%-----------------------------------------------------------------------
%  bfp_error_paper_siam.tex
%  English version of bfp_error_paper.tex, re-typeset for the SIAM class
%  siamart251216 (submission format: class + siamplain.bst + references.bib)
%
%  Compile: pdflatex paper && bibtex paper && pdflatex paper && pdflatex paper
%-----------------------------------------------------------------------
\documentclass[review,hidelinks,oneeqnum,onetabnum,onefignum,onethmnum]{siamart251216}

% The class already loads amsmath[leqno], ntheorem, cleveref, hyperref,
% xcolor, algorithm, epstopdf, xr-hyper, ifpdf and hypcap.
\usepackage{amsfonts}   % \mathbb (not loaded by the class)
\usepackage{amssymb}
\usepackage{graphicx}
\usepackage{booktabs}
\usepackage{multirow}
\usepackage{enumitem}
\usepackage{microtype}

% Allow slightly looser lines rather than overfull boxes.
\setlength{\emergencystretch}{2em}

% Compact list spacing (the manuscript has many short lists; the default
% 10pt skipped-item spacing costs several pages when summed over the paper).
\setlist{nosep,leftmargin=1.6em}
\setlist[itemize,1]{topsep=3pt,parsep=0pt,itemsep=2pt}
\setlist[enumerate,1]{topsep=3pt,parsep=0pt,itemsep=2pt}

% Tighter float separation.
\setlength{\textfloatsep}{8pt plus 2pt minus 2pt}
\setlength{\intextsep}{8pt plus 2pt minus 2pt}
\setlength{\floatsep}{8pt plus 2pt minus 2pt}

% Tighten the vertical space around displays (the paper is display-heavy).
\AtBeginDocument{%
  \setlength{\abovedisplayskip}{6pt plus 2pt minus 2pt}%
  \setlength{\belowdisplayskip}{6pt plus 2pt minus 2pt}%
  \setlength{\abovedisplayshortskip}{3pt plus 1pt}%
  \setlength{\belowdisplayshortskip}{3pt plus 1pt}%
  \setlength{\jot}{3pt}%
}

% Theorem-like environments: use the SIAM macros, not \newtheorem.
\newsiamthm{conjecture}{Conjecture}
\crefname{conjecture}{Conjecture}{Conjectures}
\newsiamremark{remark}{Remark}
% Already defined by the class: theorem, lemma, corollary, proposition,
% definition, and proof.

\headers{Rounding-Error Independence in Block Floating Point}{Y. Xiong}

\title{The Independence Assumption for Rounding Errors in Block Floating
Point: Failure of Conditional Symmetry and Its Effect on Summation Error
Accumulation\thanks{Preprint, September 2026.
\funding{The author received no external funding for this work.}}}

\author{Yongkang Xiong\thanks{College of Computer Science and Technology,
Nanjing University of Aeronautics and Astronautics, 29 Jiangjun Avenue,
Nanjing 211106, China (\email{nuaapanda@outlook.com}).}}

\begin{document}

\maketitle

\begin{abstract}
Error analyses of block floating point (BFP) and related shared-exponent formats,
including the MXFP family used by modern tensor cores, usually assume that
within-block quantization errors are zero-mean and mutually uncorrelated. That
assumption dates to a 1996 analysis of BFP roundoff and is rarely re-tested on
modern data. We re-examine it by Monte Carlo simulation and reduce the outcome to
a parameter-free model.

Four rigorous facts carry the study. Without extra assumptions,
$\operatorname{Cov}(e_i,e_j)=\mathbb{E}[\operatorname{Cov}(e_i,e_j\mid\xi)]
+\operatorname{Var}(m(\xi))$ with $m(\xi)=\mathbb{E}[e\mid\xi]$, so any failure of
the classical model is a block-level effect. The variance inflation factor
satisfies $0\le\mathcal I_n\le n$. The variance law
$\mathbb{E}[e^2]=\mathbb{E}[s^2\gamma(\xi)]$ makes the classical model exactly the
case $\gamma\equiv1/12$, $\delta\equiv0$, and separates the two defects:
$\delta\not\equiv0$ creates correlation, $\gamma\not\equiv1/12$ biases the
variance formula. Under unbounded support with $\mathbb{E}[x_1^2]<\infty$,
$m(\xi)\to-\mathbb{E}[x_1]$, $\operatorname{Var}(m(\xi))\to0$ and
$\mathcal I_n/n\to0$, so linear accumulation is transient.

The reduced model writes $m(\xi)=s\,\delta(s)$ and
$\mathbb{E}[e^2\mid\xi]=s^2\gamma(s)$: one-dimensional response functions times
the block maximum, with no adjustable parameter. Over 36 round-to-nearest
configurations it matches measured $\mathcal I_n$ within $4.5\%$ (median
$1.15\%$). Simulations cover six distributions and four mantissa widths. Zero-mean
fails for asymmetric inputs (lognormal $z=-126.8$, ReLU $z=-39.5$) and holds for
symmetric ones. For ReLU the measured variance-to-classical ratio is $0.506$,
matching the atom prediction $1-p_0=0.500$; under stochastic rounding the same
argument predicts $1.000$ against a measured $1.015$. A decomposition given $\xi$
shows that within-block correlation comes almost entirely from
$\operatorname{Var}(m(\xi))$. Bounded support forces $\mathcal I_n\to1$, whereas
unbounded support gives an inflation that first rises and then falls; under
stochastic rounding all $24$ configurations give $\mathcal I_n\approx1$.
\end{abstract}

\begin{keywords}
block floating point, low-precision arithmetic, rounding error analysis,
probabilistic error bounds, stochastic rounding, tensor cores
\end{keywords}

\begin{MSCcodes}
65G50, 65F45, 60F99, 68M15
\end{MSCcodes}

\section{Introduction}
\label{sec:intro}

Low-precision arithmetic now dominates many workloads on modern accelerators.
NVIDIA tensor cores, AMD matrix cores, and the OCP microscaling (MX) formats
raise throughput by sharing a single exponent among several elements. Formats
of this kind are called block floating point (BFP): the $n$ values of a block
share one exponent $\xi$, and each element stores only a mantissa. Relative to
elementwise floating point this saves exponent bits but reduces the dynamic
range inside the block.

For decades, error analyses of such formats have relied on a statistical model
inherited from Kallioj\"arvi and Astola \cite{kalliojarvi1996}: quantization
errors are zero-mean, their variance is determined by the block exponent and
the mantissa width, and, in later citations, the errors may also be treated as
uncorrelated \cite{lian2019,song2018,fan2018,han2025}. In parallel, numerical
analysis developed probabilistic rounding error bounds: under zero-mean and
(approximately) independent errors, summation and prefix errors grow like
$O(\sqrt n\,u)$ rather than $O(nu)$
\cite{higham2002,higham2022,hallman2023}, and the framework has been extended
to mixed precision and stochastic rounding
\cite{higham2019,elarar2024,bhola2024a,bhola2024b,croci2022,drineas2024,elarar2026},
to multiword tensor-core arithmetic \cite{fasi2023}, and to integer-splitting
emulation \cite{abdelfattah2025}.

The two communities have not fully met: hardware-oriented papers often assume the
statistical model without re-validating it on modern formats and data, while
numerical analysis papers test such models on elementwise floating point, multiword
arithmetic, integer slicing, or stochastic rounding, none of which contains a
data-determined shared exponent. This paper studies that missing case directly.

\medskip
\noindent\textbf{Contributions.}

\begin{enumerate}
\item Covariance decomposition (\cref{prop:decomp}): without further
assumptions,
$\operatorname{Cov}(e_i,e_j)=\mathbb{E}[\operatorname{Cov}(e_i,e_j\mid\xi)]
+\operatorname{Var}(m(\xi))$, so any failure of the classical model is a
block-level effect. Conditional exchangeability, usually stated as an extra
hypothesis, is automatic here (\cref{lem:exch}) and yields
$0\le\mathcal I_n\le n$ (\cref{cor:bounds}).

\item Exact variance laws and a reduced model
(\cref{prop:var}, \cref{lem:cond}): $\mathbb{E}[e^2]=\mathbb{E}[s^2\gamma(\xi)]$
and $m(\xi)=s\,\delta(s)$ with $\delta,\gamma$ depending only on the
distribution and the step, so all remaining randomness is the block maximum.
There is no adjustable parameter.

\item Exact stochastic-rounding counterpart (\cref{prop:sr}): the conditional
covariance and $\operatorname{Var}(m(\xi))$ both vanish, so
$\mathcal I_n\equiv1$ and $\operatorname{Var}(e)=\mathbb{E}[s^2\theta(1-\theta)]$
(\cref{cor:srvar}). Stochastic rounding is therefore a falsifiable probe of the
mechanism, and not merely a better rounding rule.

\item Monte Carlo experiments over six distributions and four mantissa widths:
failure of zero-mean for asymmetric inputs (\cref{sec:bias}); structural bias
of the classical variance formula with its atom explanation
(\cref{sec:var}); rise and subsequent fall of the inflation with block length
(\cref{sec:infl}); a direct decomposition given $\xi$ together with a numerical
test of (A2) (\cref{sec:cond}); and a parameter-free model that stays within
$4.5\%$ of measured $\mathcal I_n$ over $36$ configurations
(\cref{sec:model}).

\item Two support regimes (\cref{sec:domains,sec:open}): bounded support gives
$\mathcal I_n\to1$ exponentially fast and independently of the endpoint tail
index (\cref{prop:bound}), while unbounded support with
$\mathbb{E}[x_1^2]<\infty$ gives $\rho_n\to0$ and $\mathcal I_n/n\to0$
(\cref{prop:decay}). No admissible distribution in this setting yields
asymptotic linear accumulation $\mathcal I_n\asymp n$.
\end{enumerate}

\medskip
\noindent\textbf{Relation to concurrent work.} Zhang et al.\
\cite{zhang2026bfp} already observed, in the fault-tolerance context of a
BFP-based NPU, that the shared exponent couples the elements of a block; our
contribution is to quantify that coupling as a pairwise covariance identity with
separately falsifiable predictions
(\crefrange{prop:decomp}{prop:var}, \cref{prop:bound,prop:decay}). Ang et
al.\ \cite{ang2026quant} obtain
$\mathbb{E}[(S-\hat S)^2]=k\,\mathbb{E}[D^2]+k(k-1)(\mathbb{E}[D])^2$ for scalar
quantization in matrix multiplication, where the quadratic term can be ``designed
away'' by a companding quantizer; a shared exponent is the complementary case,
because $\operatorname{Var}(m(\xi))$ persists as $n$ grows, so raising the
quantization precision cannot remove the quadratic term. Sao et
al.\ \cite{sao2026trees} obtain a structurally identical centred/non-centred
dichotomy for reduction trees, but condition on the state of the tree and assume
conditionally unbiased rounding. The grid biases of Zhao et al.\
\cite{zhao2026shrinkage} and Wu et al.\ \cite{wu2026mxfp4} are defined in
\emph{amplitude space} and stay sign-symmetric for sign-symmetric inputs, so they
do not amount to $\mathbb{E}[e_i]\ne0$; our claim is that the expectation of the
signed error fails when the \emph{data} are conditionally asymmetric
(\cref{sec:bias}). Finally, \cref{prop:sr} is cognate with dither theory:
Schuchman \cite{schuchman1964} and Gray--Stockham \cite{gray1993dithered} long
ago showed that suitable dithering makes the quantization error independent of
the input; what we add is the use of that fact---driving the unobservable
$m(\xi)$ to zero---as a falsifiable probe.

\noindent This paper is predominantly a numerical study: theorem-level results
are confined to \crefrange{prop:decomp}{prop:var}, \cref{lem:cond},
\crefrange{cor:bounds}{cor:srvar}, and \cref{prop:bound,prop:decay}; all
remaining statements about scaling behaviour are explicitly labelled as numerical
observations or as unproved.

\section{Background: an inherited and never re-validated assumption}
\label{sec:background}

\subsection{Block floating point formats}

Fix a block length $n$ and a mantissa width $L_m$. The shared exponent of the
real values $x_1,\dots,x_n$ in a block is taken to be
\begin{equation}
\xi=\Bigl\lfloor \log_2 \max_{1\le i\le n}|x_i|\Bigr\rfloor ,
\qquad
s=2^{\xi-L_m},
\label{eq:format}
\end{equation}
where $s$ is called the \emph{step} of the block. The quantized value of the
$i$th element is $\hat x_i=\mathrm{RN}(x_i/s)\cdot s$, where $\mathrm{RN}$
denotes rounding to nearest, ties to even. The error is defined as
\begin{equation}
e_i=\hat x_i-x_i,\qquad i=1,\dots,n .
\label{eq:err}
\end{equation}
Equations \cref{eq:format} and \cref{eq:err} are an idealized model of the OCP
MX format family and of classical BFP: a shared exponent, rounding to nearest,
no overflow or underflow. Real hardware may use other rounding rules or
sticky-bit handling; \cref{sec:modes} analyses the sensitivity to the rounding
rule, and a comparison against real hardware is left to future work
(\cref{sec:limitations}).

\subsection{The inherited assumption and its propagation chain}

In a deterministic signal-processing setting, Kallioj\"arvi and Astola
\cite{kalliojarvi1996} established two properties of the BFP quantization error,
which we call the \emph{classical model}:
\begin{equation}
\text{(H1, zero mean)}\quad \mathbb{E}[e_i]=0,\qquad
\text{(H2, variance)}\quad \operatorname{Var}[e_i]=\mathbb{E}[s^2]/12 .
\label{eq:H1H2}
\end{equation}
Later work inherited (H1) and (H2) when citing that result, and also introduced a
third, stronger and much less discussed assumption,
\begin{equation}
\text{(H3, uncorrelated)}\qquad \operatorname{Cov}(e_i,e_j)=0,\qquad i\ne j .
\label{eq:H3}
\end{equation}
\Cref{tab:chain} lists, for three representative papers, the property that each
one takes over from the 1996 result and how it is taken over. In every case a
property of the BFP quantization error is inherited from that result without
being re-derived under modern formats or neural-network activation
distributions.

\begin{table}[!tb]
\centering
\caption{The propagation chain of the classical assumption. For each reference
the table states the property inherited from the 1996 result and how.}
\label{tab:chain}
\setlength{\tabcolsep}{5pt}
\small
\begin{tabular}{@{}p{0.28\textwidth}p{0.64\textwidth}@{}}
\toprule
Reference & Inherited property \\
\midrule
Song, Liu, and Wang, AAAI 2018 \cite{song2018} &
take the quantization error of a block to have zero mean and variance
$\sigma^2$, attributing both directly to Kallioj\"arvi and Astola (1996) \\
\addlinespace
Lian et al., IEEE TVLSI 2019 \cite{lian2019} &
model $e_{bm}$, the quantization error of the block mantissa, as an uncorrelated
random variable rather than deriving zero-meanness for a data-determined
shared exponent \\
\addlinespace
Han et al., DAC 2025 \cite{han2025} &
state that for a block floating point using round to nearest the quantisation
error is zero-mean with variance $\sigma^2$, citing their [31], which is the 1996
reference \\
\bottomrule
\end{tabular}
\end{table}

The same inheritance appears for modern MX formats. Wu et al.\
\cite{wu2026mxfp4} note that ``existing work treats the quantization error as a
monolithic noise term'', while their own three-way decomposition of the MXFP4
error (scale bias, dead zone, grid noise) is defined in amplitude space and does
not touch the conditional mean of the signed error. One citation deserves a
qualification: Fan et al.\ \cite{fan2018} do cite the 1996 reference, but only in
the historical sense that BFP has long been used in FPGA designs, and they
attribute the CNN error analysis instead to Song, Liu, and Wang; we therefore
exclude them from \cref{tab:chain}.

The original application context of the 1996 work matters here: audio, radar, and
array signal processing, whose inputs are approximately symmetric within a block.
Symmetry is therefore the substantive condition for (H1) and (H3), whereas modern
neural-network activations---especially post-ReLU activations---are nonnegative,
so conditional symmetry is absent.

\subsection{Progress on the numerical analysis side}

The foundations of rigorous probabilistic error analysis were laid by Higham
and Mary and co-workers \cite{higham2019}; subsequent developments include
probabilistic summation bounds \cite{hallman2023}, deterministic and
probabilistic bounds for mixed-precision algorithms
\cite{bhola2024a,bhola2024b}, probabilistic analyses of stochastic rounding
\cite{elarar2024,drineas2024,elarar2026}, low-precision PDE solution
\cite{croci2022}, multiword tensor-core arithmetic and integer-splitting matrix
multiplication \cite{fasi2023,abdelfattah2025}, mixed-precision iterative
refinement \cite{carson2024}, BFP acceleration of CNNs \cite{fan2018}, and
tooling for low-precision emulation \cite{carson2025}.

All of these works rely on (H1) and on independence, or at least on
uncorrelatedness; but none of the objects analysed contains a data-determined
shared exponent. In \cref{eq:format} the exponent $\xi$ is a function of all
elements of the block, which makes the within-block errors intrinsically
coupled through $\xi$ (contrast Sao et al.\ \cite{sao2026trees} and
Fox--Lindstrom \cite{fox2024zfp}).

\section{Formalization and rigorous results}
\label{sec:formal}

\subsection{Notation}

Let the elements of a block be i.i.d.\ with finite error second moments, and let
\begin{equation}
m(\xi)=\mathbb{E}[e_1\mid\xi],\quad
c(\xi)=\operatorname{Cov}(e_i,e_j\mid\xi)\ (i\ne j),\quad
\rho=\frac{\operatorname{Cov}(e_i,e_j)}{\operatorname{Var}(e)}\ (i\ne j).
\label{eq:notation}
\end{equation}

\begin{lemma}[Conditional exchangeability holds automatically]
\label{lem:exch}
If the block elements $x_1,\dots,x_n$ are i.i.d.\ and $\xi$ is a symmetric
function of $(|x_1|,\dots,|x_n|)$, then $(e_1,\dots,e_n)$ is exchangeable given
$\xi$; consequently $m(\xi)$ and $c(\xi)$ do not depend on the indices, and
$\rho$ does not depend on which pair $i\ne j$ is chosen.
\end{lemma}

\begin{proof}
For any permutation $\pi$, $(x_{\pi(1)},\dots,x_{\pi(n)})$ has the same
distribution as $(x_1,\dots,x_n)$, and $\xi$ in \cref{eq:format} is invariant
under permutations, so both have the same conditional law given $\xi$. Since
$e_i$ depends on the data only through $(x_i,\xi)$, the vectors
$(e_{\pi(1)},\dots,e_{\pi(n)})$ and $(e_1,\dots,e_n)$ are identically distributed
given $\xi$.
\end{proof}

The variance of the within-block error sum is
\begin{equation}
\operatorname{Var}\Bigl[\sum_{i=1}^{n}e_i\Bigr]
= n\operatorname{Var}(e)\bigl(\underbrace{1+(n-1)\rho}_{=:\ \mathcal I_n}\bigr).
\label{eq:inflation}
\end{equation}
We call $\mathcal I_n$ the \emph{variance inflation factor}, the central quantity
of this paper: the classical model (H1)--(H3) implies $\mathcal I_n=1$ (and
conversely, only $\operatorname{Cov}(e_i,e_j)=0$ is needed, not zero means).
Conditional exchangeability, often listed as an additional hypothesis, follows
from i.i.d.\ inputs and symmetry of $\xi$ by \cref{lem:exch}, so the only
hypothesis left to test is \emph{conditional uncorrelatedness} (A2), measured in
\cref{sec:cond}.

\subsection{Covariance decomposition}

\begin{proposition}[Covariance decomposition]
\label{prop:decomp}
Let the within-block errors $(e_1,\dots,e_n)$ be exchangeable given $\xi$ and
have finite second moments. Then, \emph{without any additional assumption},
\begin{equation}
\operatorname{Cov}(e_i,e_j)=\mathbb{E}\bigl[c(\xi)\bigr]+\operatorname{Var}\bigl(m(\xi)\bigr),
\qquad i\ne j .
\label{eq:dec-exact}
\end{equation}
If in addition the conditional uncorrelatedness assumption
\begin{equation}
\textup{(A2)}\qquad
\operatorname{Cov}(e_i,e_j\mid\xi)=0,\quad i\ne j,
\quad\text{i.e.}\quad c(\xi)\equiv0,
\label{eq:A2}
\end{equation}
holds, then \cref{eq:dec-exact} reduces to
$\operatorname{Cov}(e_i,e_j)=\operatorname{Var}(m(\xi))$, and
\begin{equation}
\rho=\frac{\operatorname{Var}(m(\xi))}{\operatorname{Var}(e)},
\qquad
\mathcal I_n=1+(n-1)\frac{\operatorname{Var}(m(\xi))}{\operatorname{Var}(e)} .
\label{eq:rho}
\end{equation}
\end{proposition}

\begin{proof}
By the law of total covariance,
$\operatorname{Cov}(e_i,e_j)
=\mathbb{E}\bigl[\operatorname{Cov}(e_i,e_j\mid\xi)\bigr]
+\operatorname{Cov}\bigl(\mathbb{E}[e_i\mid\xi],\mathbb{E}[e_j\mid\xi]\bigr)$.
The first term is $\mathbb{E}[c(\xi)]$; since
$\mathbb{E}[e_i\mid\xi]=\mathbb{E}[e_j\mid\xi]=m(\xi)$ by \cref{lem:exch}, the
second is $\operatorname{Cov}(m(\xi),m(\xi))=\operatorname{Var}(m(\xi))$.
\Cref{eq:rho} follows from the definition of $\rho$ and \cref{eq:inflation}.
\end{proof}

\begin{corollary}[Rigorous bounds on the inflation factor]
\label{cor:bounds}
Always,
\begin{equation}
0\le \mathcal I_n\le n ,
\qquad\text{that is,}\qquad
\operatorname{Var}\Bigl[\sum_{i=1}^{n}e_i\Bigr]\le n^2\operatorname{Var}(e).
\label{eq:bounds}
\end{equation}
\end{corollary}

\begin{proof}
The left inequality is
$\operatorname{Var}[\sum_i e_i]=n\operatorname{Var}(e)+n(n-1)\operatorname{Cov}(e_i,e_j)\ge0$,
equivalent to $\rho\ge-1/(n-1)$. The right one follows from Cauchy--Schwarz:
$e_i$ and $e_j$ are identically distributed, so
$\operatorname{Cov}(e_i,e_j)\le\sqrt{\operatorname{Var}(e_i)\operatorname{Var}(e_j)}=\operatorname{Var}(e)$,
hence $\rho\le1$.
\end{proof}

\begin{remark}
The bound $\mathcal I_n\le n$ is substantive: the within-block error sum has
standard deviation growing at most linearly in $n$, so the \emph{asymptotic}
slope of $\log\mathcal I_n$ against $\log n$ is at most $1$ and the local slope
$1.134$ measured in \cref{sec:infl} is necessarily a transient
(\cref{sec:open} sharpens this to $\mathcal I_n/n\to0$). It also flags
impossible outputs: $\mathcal I_n<0$ or $\mathcal I_n>n$ must be wrong.
Conversely, \cref{prop:decomp} attributes the within-block correlation entirely
to the fluctuation of $m(\xi)$ and the mean of $c(\xi)$: if $m(\xi)\equiv0$ and
$c(\xi)\equiv0$---which \cref{sec:sr} proves for stochastic rounding---then (H3)
holds exactly and $\mathcal I_n\equiv1$ however large $\operatorname{Var}(e)$ may
be, whereas a nondegenerate $m(\xi)$ gives \emph{every} pair the same nonzero
covariance, accumulating like $n^2$ rather than $n$ in a summation.
\end{remark}

Using $e_i=s\cdot\bigl(\mathrm{RN}(x_i/s)-x_i/s\bigr)$, the conditional mean is
\begin{equation}
m(\xi)=s\cdot\delta(\xi),\qquad
\delta(\xi)=\mathbb{E}\bigl[\mathrm{RN}(u)-u \mid \xi\bigr],\quad u=x/s,
\label{eq:mxi}
\end{equation}
where $\delta(\xi)$ is the mean unit-step rounding residual given $\xi$;
$\delta(\xi)=0$ whenever the conditional distribution of $u$ is symmetric about
$0$.

\subsection{An exact counterpart for stochastic rounding}
\label{sec:srtheory}

\begin{proposition}[Zero inflation under stochastic rounding]
\label{prop:sr}
Let the quantization use stochastic rounding (SR): given $x$ and $s$, put
$u=x/s$, $k=\lfloor u\rfloor$ and $\theta=u-k\in[0,1)$, and take the values
$k$ and $k+1$ with probabilities $1-\theta$ and $\theta$, respectively, with
the randomizations mutually independent given $(x_1,\dots,x_n)$. Then
$\mathbb{E}[e_i\mid x,\xi]=0$, $m(\xi)\equiv0$,
$\operatorname{Cov}(e_i,e_j)=0$ $(i\ne j)$ and $\mathcal I_n\equiv1$.
\end{proposition}

\begin{proof}
Given $x$ and $s$, the definition of SR gives
$\mathbb{E}[\mathrm{RN}(u)]=k(1-\theta)+(k+1)\theta=k+\theta=u$, hence
$\mathbb{E}[e\mid x,s]=s\bigl(\mathbb{E}[\mathrm{RN}(u)]-u\bigr)=0$ and, by
conditional expectation,
$m(\xi)=\mathbb{E}\bigl[\mathbb{E}[e\mid x,\xi]\mid\xi\bigr]=0$. Since the SR
randomizations of the individual elements are mutually independent given
$(x_1,\dots,x_n)$, $\operatorname{Cov}(e_i,e_j\mid\xi,x)=0$; the law of total
covariance then gives $\operatorname{Cov}(e_i,e_j\mid\xi)=0$, because both
$\mathbb{E}[e_i\mid x,\xi]$ and $\mathbb{E}[e_j\mid x,\xi]$ vanish. Hence (A2)
holds automatically for SR, and \cref{prop:decomp} gives
$\operatorname{Cov}(e_i,e_j)=\operatorname{Var}(m(\xi))=0$.
\end{proof}

\begin{corollary}[Variance law under stochastic rounding]
\label{cor:srvar}
Write $\theta=\operatorname{frac}(x/s)\in[0,1)$. Under stochastic rounding,
\begin{equation}
\mathbb{E}[e^2\mid\xi]=s^2\,\mathbb{E}\bigl[\theta(1-\theta)\mid\xi\bigr],
\qquad
\operatorname{Var}(e)=\mathbb{E}\bigl[s^2\theta(1-\theta)\bigr].
\label{eq:srvar}
\end{equation}
In particular, if $\theta$ is conditionally uniform on $[0,1)$ given $\xi$, then
$\operatorname{Var}(e)=\mathbb{E}[s^2]/6$, \emph{exactly twice the classical
round-to-nearest value $\mathbb{E}[s^2]/12$}.
\end{corollary}

\begin{proof}
Given $x$ (hence $u$ and $\theta$), $\mathrm{RN}(u)$ takes the values $k,k+1$ with
probabilities $1-\theta,\theta$, so $g=\mathrm{RN}(u)-u$ equals $-\theta$ with
probability $1-\theta$ and $1-\theta$ with probability $\theta$, giving
$\mathbb{E}[g^2\mid x,s]=\theta(1-\theta)$ and $\mathbb{E}[g\mid x,s]=0$.
Multiplying by $s^2$ and taking conditional expectations gives the first
identity, and the second follows from $\operatorname{Var}(e)=\mathbb{E}[e^2]$
since $m\equiv0$. If $\theta\mid\xi$ is uniform,
$\mathbb{E}[\theta(1-\theta)]=1/6$, so $\operatorname{Var}(e)=\mathbb{E}[s^2]/6$.
\end{proof}

\Cref{prop:sr} provides an operational probe: in a shared-exponent block,
stochastic rounding drives the unobservable $m(\xi)$ to zero, so if the
inflation in \cref{sec:infl} is caused by $m(\xi)\ne0$, switching to SR must
annihilate it exactly, independently of the distribution, block length, and
mantissa width. \Cref{sec:sr} carries out this test.

\subsection{Exact variance laws and the reduced model}
\label{sec:modeltheory}

Both $m(\xi)$ and the variance structure reduce to one-dimensional functions,
turning the scaling question into a computable object. With
$g(u)=\mathrm{RN}(u)-u\in[-\tfrac12,\tfrac12]$, $e_i=s\,g(x_i/s)$.

\begin{proposition}[Variance law]
\label{prop:var}
Write $\gamma(\xi)=\mathbb{E}\bigl[g(x_1/s)^2\mid\xi\bigr]$ and
$\delta(\xi)=\mathbb{E}\bigl[g(x_1/s)\mid\xi\bigr]=m(\xi)/s$. Then
\begin{equation}
\mathbb{E}[e^2]=\mathbb{E}\bigl[s^2\gamma(\xi)\bigr],
\qquad
\operatorname{Var}(e)=\mathbb{E}\bigl[s^2\gamma(\xi)\bigr]-\bigl(\mathbb{E}[e]\bigr)^2 .
\label{eq:vardecomp}
\end{equation}
The classical model (H1)--(H2) is exactly equivalent to $\delta\equiv0$ and
$\gamma\equiv1/12$.
\end{proposition}

\begin{proof}
Taking expectations in $e^2=s^2g(x/s)^2$ (with $s$ a function of $\xi$) gives
the first identity; the second follows from
$\operatorname{Var}(e)=\mathbb{E}[e^2]-(\mathbb{E}[e])^2$. Moreover
$\mathbb{E}[e^2]=\mathbb{E}[s^2/12]$ if and only if $\gamma\equiv1/12$, and
$\mathbb{E}[e]=0$ if and only if $\delta\equiv0$.
\end{proof}

\Cref{prop:var} shows that the failure of the classical variance and the
emergence of correlation are \emph{two independent structural defects}, the
former $\gamma\not\equiv1/12$ and the latter $\delta\not\equiv0$, so the
deviations of \cref{sec:var} and \cref{sec:infl} cannot mask each other.

The following lemma shows that, given $\xi$, the conditional distribution enters
the behavior of $g$ only through the step $s$, with error controlled by the
distance of the block maximum from the bulk; this legitimizes the reduced model
below.

\begin{lemma}[Two-piece mixture representation of the conditional law]
\label{lem:cond}
Write $a=2^{\xi}$, $b=2^{\xi+1}$, $F_a=F(a)$, $F_b=F(b)$ and
$P_k=\mathbb P(\xi=k)=F_b^n-F_a^n$. Given $\xi=k$, the conditional law of $x_1$ is
a mixture of two truncated distributions,
\begin{equation}
\mathbb P(x_1\in A\mid\xi=k)
=w_A\,\mathbb P_F\bigl(x_1\in A\mid |x_1|<a\bigr)
+w_B\,\mathbb P_F\bigl(x_1\in A\mid a\le|x_1|<b\bigr),
\label{eq:conddens}
\end{equation}
with
$w_A=F_a\bigl(F_b^{n-1}-F_a^{n-1}\bigr)/P_k$,
$w_B=F_b^{n-1}\bigl(F_b-F_a\bigr)/P_k$ and $w_A+w_B=1$. Consequently, for every
bounded measurable $\varphi$,
\begin{equation}
\bigl|\mathbb E[\varphi(x_1)\mid\xi=k]-\mathbb E_F[\varphi]\bigr|
\;\le\;2\|\varphi\|_\infty\bigl[w_A(1-F_a)+w_B(1-F_b+F_a)\bigr].
\label{eq:condbound}
\end{equation}
If the support of $F$ is unbounded, then as $k\to\infty$ we have $F_a,F_b\to1$,
$w_A\to(n-1)/n$ and $w_B\to1/n$, so the right-hand side of \cref{eq:condbound}
tends to $0$.
\end{lemma}

\begin{proof}
Given $x_1=x$ with $|x|<a$, the event $\{\xi=k\}$ is equivalent to
$\max_{j\ge2}|x_j|\in[a,b)$, of probability $F_b^{n-1}-F_a^{n-1}$; given
$a\le|x|<b$ it is equivalent to $\max_{j\ge2}|x_j|<b$, of probability
$F_b^{n-1}$. Multiplying these by $f(x)/P_k$ yields \cref{eq:conddens}, whose
normalization follows from
$F_a(F_b^{n-1}-F_a^{n-1})+F_b^{n-1}(F_b-F_a)=F_b^n-F_a^n$. The error bound follows
from
$\mathbb E[\varphi\mid\xi=k]-\mathbb E_F[\varphi]=\sum_{i\in\{A,B\}}w_i\bigl(\mathbb E_F[\varphi\mid A_i]-\mathbb E_F[\varphi]\bigr)$
together with
$|\mathbb E_F[\varphi\mid A_i]-\mathbb E_F[\varphi]|\le2\|\varphi\|_\infty\bigl(1-\mathbb P_F(A_i)\bigr)$,
where $1-\mathbb P_F(A)=1-F_a$ and $1-\mathbb P_F(B)=1-F_b+F_a$.
\end{proof}

\begin{remark}
\Cref{eq:condbound} makes precise when ``the conditional law enters only through
the step'' is valid: the deviation is controlled by two \emph{truncations}, at
$a$ and at $b$. It degenerates when $2^\xi$ lies far below the bulk scale of the
data ($F_a\to0$)---substantive information, not a technical defect: almost all
elements then fall inside the rounding dead zone and the conditional law
genuinely differs from the unconditional one, its response governed by the other
asymptotics \cref{eq:lmscaling}. Conversely, as $2^\xi$ grows the truncations
disappear, so the reduced model of \cref{sec:model} is more accurate on the
\emph{large}-$\xi$ side and \cref{prop:decay} rests on the same limit. Where the
right-hand side of \cref{eq:condbound} is negligible, $\varphi(x)=g(x/s)$ (with
$\|\varphi\|_\infty=\tfrac12$) gives
\begin{equation}
\begin{gathered}
\delta(\xi)=\delta(s)+O(\varepsilon),\qquad
\gamma(\xi)=\gamma(s)+O(\varepsilon),\\[2pt]
\delta(s)=\mathbb E_F\bigl[\mathrm{RN}(x/s)-x/s\bigr],\qquad
\gamma(s)=\mathbb E_F\bigl[\bigl(\mathrm{RN}(x/s)-x/s\bigr)^2\bigr],
\end{gathered}
\label{eq:deltagamma}
\end{equation}
with $\varepsilon$ that right-hand side; over the experimental regime this
remainder has size $10^{-3}$, which is why the reduced model reproduces the
simulations to three or four significant digits.
\end{remark}

\noindent\textbf{Reduced model.} Retaining the leading term of
\cref{eq:deltagamma} yields the central claim: every statistic of the
within-block error is determined by two one-dimensional functions of the
distribution and the step,
\begin{equation}
m(\xi)=s\,\delta(s),\qquad
\mathbb{E}[e^2\mid\xi]=s^2\gamma(s),
\qquad s=2^{\xi-L_m},
\label{eq:reduced}
\end{equation}
so that
\begin{equation}
\begin{gathered}
\operatorname{Var}(e)=\mathbb{E}\bigl[s^2\gamma(s)\bigr]-\bigl(\mathbb{E}[s\,\delta(s)]\bigr)^2,
\qquad
\operatorname{Cov}(e_i,e_j)=\operatorname{Var}\bigl(s\,\delta(s)\bigr),\\[2pt]
\rho=\frac{\operatorname{Var}(s\,\delta(s))}{\operatorname{Var}(e)},
\qquad
\mathcal I_n=1+(n-1)\frac{\operatorname{Var}(s\,\delta(s))}{\operatorname{Var}(e)} .
\end{gathered}
\label{eq:reduced2}
\end{equation}
In \cref{eq:reduced,eq:reduced2} the only random input is the shared step $s$
(the rounded logarithmic scale of $M_n$), while $\delta$ and $\gamma$ follow
from one-dimensional quadrature: the problem factorizes into a
\emph{deterministic response curve} and an \emph{order statistic} (the
distribution of $M_n$). Three consequences follow.
\begin{itemize}
\item The symmetry criterion is exact: if $F$ is symmetric about $0$, then
$g(\cdot/s)$ is odd, so $\delta(s)\equiv0$, hence
$\operatorname{Cov}(e_i,e_j)=0$ and $\mathcal I_n\equiv1$ exactly, independently
of (A2).

\item The scaling question reduces to an extreme-value question: the growth rate
of $\mathcal I_n-1$ is governed entirely by whether
$B_n=\operatorname{Var}(s\,\delta(s))$ decays with $n$, and the law of $s$
depends on $n$ only through $M_n$; with \cref{cor:bounds} this bounds the
exponent by $1$, and \cref{prop:decay} in \cref{sec:open} yields $B_n\to0$.

\item The $L_m$ dependence has a quantitative explanation. In the ``dead zone''
regime (step below the bulk scale of the data) most elements of a block satisfy
$|x_i|<s/2$ and round to $0$; a direct computation gives
$\delta(s)=-\alpha s+o(s)$ and $\gamma(s)=\beta+o(1)$, with $\alpha,\beta>0$
depending only on the shape of $F$ near $0$. Then $m=s\delta\propto s^2$ and
$B_n\propto\operatorname{Var}(s^2)$, whereas
$\operatorname{Var}(e)\approx\mathbb{E}[s^2\gamma]\propto\mathbb{E}[s^2]$; since
one more bit of $L_m$ halves $s$,
\begin{equation}
\operatorname{Var}(s^2)\to\tfrac1{16}\operatorname{Var}(s^2),\qquad
\mathbb{E}[s^2]\to\tfrac14\mathbb{E}[s^2],\qquad
\rho\ \propto\ 4^{-L_m},\qquad \mathcal I_n-1\ \propto\ 4^{-L_m}.
\label{eq:lmscaling}
\end{equation}
The data of \cref{sec:model} support this scaling ($16.3$ for ReLU, $16.9$ for
the half-normal, close to $4^{2}=16$); it fails near the peak of $|\delta(s)|$,
which explains the non-monotone $L_m$ dependence in \cref{sec:infl}.
\end{itemize}

\subsection{The two regimes of the support}
\label{sec:domains}

The first regime is \cref{prop:bound}; the second is \cref{prop:decay} in
\cref{sec:open}, where the inflation is transient.

\begin{proposition}[Bounded support: exponential decay]
\label{prop:bound}
Suppose $|x_1|$ has bounded support, $B=\operatorname{ess\,sup}|x_1|<\infty$,
and $\mathbb{E}\bigl[\operatorname{Var}(e\mid\xi)\bigr]>0$. Let $\xi_0$ be the
integer with $2^{\xi_0}<B\le2^{\xi_0+1}$ and write $p=F(2^{\xi_0})<1$. Then
$\xi\le\xi_0$ almost surely, and
\begin{equation}
\begin{gathered}
\mathbb{P}\bigl(\xi\ne\xi_0\bigr)\le p^{\,n},
\qquad
\operatorname{Var}\bigl(m(\xi)\bigr)\le C\,p^{\,n},\\
\mathcal I_n-1\;\le\;\frac{(n-1)\,C\,p^{\,n}}{\mathbb{E}\bigl[\operatorname{Var}(e\mid\xi)\bigr]\bigl(1-p^{\,n}\bigr)}\;\longrightarrow\;0 ,
\end{gathered}
\label{eq:bounded}
\end{equation}
where $C$ depends only on $F$ and $L_m$, and the convergence is exponentially
fast.
\end{proposition}

\begin{proof}
Since $M_n\le B\le2^{\xi_0+1}$, we have $\xi\le\xi_0$ except on the
null event $\{M_n=B=2^{\xi_0+1}\}$. Moreover
$\{\xi\ne\xi_0\}\subseteq\{\xi\le\xi_0-1\}\subseteq\{M_n\le2^{\xi_0}\}$, so
$\mathbb{P}(\xi\ne\xi_0)\le F(2^{\xi_0})^n=p^n$. From $|g|\le1/2$ we get
$|m(\xi)|\le\tfrac122^{\xi_0-L_m}$, hence
$\operatorname{Var}(m)\le(\max m-\min m)^2\,\mathbb{P}(\xi\ne\xi_0)\le Cp^n$.
Finally $\operatorname{Var}(e)=\mathbb{E}[\operatorname{Var}(e\mid\xi)]+\operatorname{Var}(m)$,
and substitution into \cref{eq:rho} gives the claim.
\end{proof}

\begin{remark}
\Cref{prop:bound} is meaningful only in the nondegenerate case
$\mathbb{E}[\operatorname{Var}(e\mid\xi)]>0$. The conclusion is \emph{independent
of the endpoint tail index}: bounded support forces $\xi$ onto a single grid
point, so the inflation must vanish exponentially; heavy tails accompany
unbounded support and are thus not necessary for the inflation. The quantity
$p=F(2^{\xi_0})$ in \cref{eq:bounded} is often close to $1$ for continuous $F$,
making the bound conservative: the real mechanism is concentration of $\xi$ on
the \emph{mode} grid point (\cref{sec:domains-res} reports numerically that
$\mathbb{P}(\xi=\text{mode})$ already reaches $0.97$--$1.00$ at $n=4096$).
\end{remark}

\subsection{Unbounded support: the transient nature of the inflation}
\label{sec:open}

With unbounded support the inflation also disappears, but so slowly that over
attainable block lengths it looks like a power law.

\begin{proposition}[Asymptotic disappearance of the inflation under unbounded support]
\label{prop:decay}
Suppose $|x_1|$ has unbounded support, $\mathbb E[x_1^2]<\infty$ and
$\operatorname{Var}(x_1)>0$. Write
$A_n=\mathbb E\bigl[\operatorname{Var}(e\mid\xi)\bigr]$ and
$B_n=\operatorname{Var}\bigl(m(\xi)\bigr)$, so that
$\operatorname{Var}(e)=A_n+B_n$. Then, as $n\to\infty$,
\begin{equation}
m(\xi)\longrightarrow-\mathbb E[x_1]\quad\text{a.s.},\qquad
B_n\longrightarrow0,\qquad
A_n\longrightarrow\operatorname{Var}(x_1)>0 ,
\label{eq:decay}
\end{equation}
and hence, under assumption (A2),
\begin{equation}
\rho_n=\frac{B_n}{A_n+B_n}\longrightarrow0,
\qquad
\mathcal I_n=1+(n-1)\rho_n=\frac{A_n+nB_n}{A_n+B_n},
\qquad
\frac{\mathcal I_n}{n}\le\frac1n+\frac{B_n}{A_n}\longrightarrow0 .
\label{eq:decay2}
\end{equation}
In particular, no unbounded distribution with finite second moment yields
$\mathcal I_n\asymp n$: linear accumulation is at best a transient over moderate
block lengths.
\end{proposition}

\begin{proof}
Write $s=2^{\xi-L_m}$ and $D=\{|x_1|<s/2\}$ (the dead zone). By \cref{lem:cond},
given $\xi=k$ the conditional law of $x_1$ is a mixture of the truncated
components $A=\{|x_1|<2^k\}$ and $B=\{2^k\le|x_1|<2^{k+1}\}$ with weights
$w_A\to(n-1)/n$ and $w_B\to1/n$ as $k\to\infty$.

\emph{Component $B$.} On it $|x_1|\ge2^k=s2^{L_m}$, so
$w_B\le\mathbb E\bigl[|x_1|;2^k\le|x_1|<2^{k+1}\bigr]/2^k\to0$ (because
$\mathbb E|x_1|<\infty$); and on it $|e_1|\le s/2$, so the component contributes
at most
$s w_B/2\le\mathbb E\bigl[|x_1|;|x_1|\ge2^k\bigr]/2^{L_m+1}\to0$.

\emph{Component $A$.} There
$\mathbb P(D^c\mid A)\le\bar F(s/2)/F(2^k)\to0$, and on $D^c$ we have
$|e_1|\le s/2$, so the contribution is at most
$(s/2)\bar F(s/2)/F(2^k)\to0$ (because $t\bar F(t)\to0$); on $D$, $e_1=-x_1$
with truncation bound $s/2\to\infty$, so
$\mathbb E[e_1\mathbb 1_D\mid A]\to-\mathbb E[x_1]$.

Combining the components gives $m(\xi=k)\to-\mathbb E[x_1]$ (the exact limit
is $-(1-n^{-1})\mathbb E[x_1]$). By the estimates above,
$|m(\xi)|\le w_A\mathbb E[|x_1|\mid A]+\tfrac s2w_B\le2\mathbb E|x_1|$ is
uniformly bounded, so dominated convergence gives
$\mathbb E[m(\xi)]\to-\mathbb E[x_1]$ and
$\mathbb E[m(\xi)^2]\to(\mathbb E[x_1])^2$, that is, $B_n\to0$. Replacing
$\varphi$ by $g(x/s)^2$ and repeating the argument gives
$A_n\to\operatorname{Var}(x_1)$. Finally, the law of total variance with
$\rho_n=B_n/(A_n+B_n)$ from \cref{prop:decomp}, substituted into
\cref{eq:inflation}, yields \cref{eq:decay2}.
\end{proof}

\begin{remark}[Peak and fall-off: numerical evidence from the model]
\Cref{prop:decay} gives only the limit, not the rate. Evaluating the right-hand
side of \cref{eq:reduced2} \emph{exactly} by one-dimensional quadrature (not
Monte Carlo; Pareto$(4)$, $L_m=2$, $60$ decimal digits to eliminate cancellation
error) yields the whole population curve of $\mathcal I_n$: it peaks at
$5.9\times10^4$ near $n\approx10^5$, then falls monotonically back to $1$ ($52$
at $n=10^{12}$, $1.002$ at $n=10^{20}$), and reproduces the measured values
(\cref{tab:bign}): $3.61\times10^4$ at $n=10^6$ and $9.2\times10^3$ at $n=10^7$,
against measured $3.61\times10^4$ and $9.6\times10^3$. The fall-off is therefore
\emph{population} behavior, not a sampling artifact (see \cref{sec:inflbig}).
The dependence on $n$ is also not smooth: rounding in
$\xi=\lfloor\log_2M_n\rfloor$ makes $B_n$ oscillate with $n$, so a power-law fit
over a finite range is inappropriate.
\end{remark}

\begin{remark}[What remains open]
\Cref{eq:decay2} gives $\mathcal I_n=o(n)$ but not $\mathcal I_n\to1$: the latter
is equivalent to $nB_n\to0$ and would require an \emph{$n$-independent}
regularity assumption on $\delta(s)$ as $s\to\infty$, which we do not state here.
Moreover, for $\alpha\le4$ the fourth moment of $e$ diverges, so the sample
estimator of $\mathcal I_n$ has no finite variance; the limit \cref{eq:decay} is
therefore accessible only by high-precision evaluation of the model, not
``verified'' by simulation with a fixed number of blocks.
\end{remark}

\section{Numerical experiment design}
\label{sec:methods}

\subsection{Distributions and format parameters}

The experiments use the six distributions listed in \cref{tab:dists}, chosen to
separate the two factors ``symmetry'' and ``heavy tails'':

\begin{table}[!tb]
\centering
\caption{The six distributions used in the experiments.}
\label{tab:dists}
\small
\begin{tabular}{@{}llp{0.44\textwidth}@{}}
\toprule
Distribution & Symmetry & Support and rationale \\
\midrule
Gaussian & symmetric & unbounded; baseline and control \\
Student-$t_3$ & symmetric & unbounded; heavy-tailed control (separates heavy tails from symmetry)\\
Spike mixture & symmetric & unbounded; outlier channel $95\%\ \mathcal N(0,0.01^2)+5\%\ \mathcal N(0,1)$ \\
Lognormal & asymmetric & unbounded; strongly asymmetric, heavy right tail \\
Half-normal & asymmetric & bounded on one side; intermediate case \\
ReLU output & asymmetric & bounded on one side; the standard case for neural-network activations \\
\bottomrule
\end{tabular}
\end{table}

The mantissa widths are $L_m=2,3,4,7$, corresponding to E5M2, E4M3/MXFP4,
4-bit and bf16 scale, respectively. The shared exponent and the step are
computed as in \cref{eq:format}, and rounding is round-to-nearest (ties to even)
unless stated otherwise. All experiments use fixed random seeds. The
reproduction scripts, the raw outputs of the runs reported here, and a table
mapping each script to the results it produces are openly available at
\url{https://github.com/NUAApanda-1116/bfp-error-independence}.

\subsection{The estimator and the sequence of corrections}
\label{sec:estimator}

By \cref{eq:inflation}, $\mathcal I_n$ can be recovered from $\rho$ or estimated
directly. Our estimator went through three rounds of correction, recorded here
because the first two are instructive errors.

\paragraph{Two rejected rounds} The initial implementation computed the within-block
correlation after centring within each block, which forces the within-block error
sum to zero and hence makes the average pairwise correlation \emph{necessarily}
equal to $-1/(n-1)$: across $48$ configurations the estimate depended only on the
block size ($-1/15$, $-1/31$, $-1/63$) and not at all on the distribution. The
second version estimated the average covariance from a single pair $(e_0,e_1)$ and
formed $1+(n-1)\rho$; its standard error is inflated by the factor $(n-1)$, and it
produced \emph{negative} inflation factors at $n=256$ and $n=512$ ($-17.487$ and
$-6.797$), whereas $\mathcal I_n\ge0$ always holds. We therefore withdraw all
scaling conclusions that rested on it.

\paragraph{The estimator used here} Rather than estimating $\rho$, we estimate the
quantity in \cref{eq:inflation} directly,
\begin{equation}
\widehat{\mathcal I_n}
=\frac{\widehat{\operatorname{Var}}_{\text{blocks}}\bigl(\sum_{i\in\text{block}}e_i\bigr)}
{n\cdot\widehat{\operatorname{Var}}_{\text{elem}}(e)} ,
\label{eq:estimator}
\end{equation}
whose numerator and denominator are both low-variance statistics, using all block
sums and all elements respectively and depending on no single pair. Error bars are
standard errors over $20$ independent groupings of $1000$ blocks each; all data in
\cref{sec:results} come from this estimator and satisfy $\mathcal I_n\ge0$.

\paragraph{A fourth tool} To measure the two terms of \cref{prop:decomp} separately
(\cref{sec:cond,sec:model}) a further script
(\texttt{bfp\_error\_pilot\_v7.py}) bins by the integer value of $\xi$, centres by
the \emph{within-bin (conditional) mean} $m_b$, and retains only
$(S,\sum_i e_i^2,\xi,s)$ per block, which by the identity
$\sum_{i\ne j}(e_i-m_b)(e_j-m_b)=(S-nm_b)^2-\sum_i(e_i-m_b)^2$ (with
$S=\sum_i e_i$) determines $c(\xi)$ exactly. The centring is essential: centring
by the within-block mean makes the average pairwise covariance necessarily equal to
$-\operatorname{Var}(e)/(n-1)$, the artefact of the first round, whereas centring
by the within-bin mean estimates $\operatorname{Cov}(e_i,e_j\mid\xi)$ directly. The
same script yields $\mathcal I_n$ along three independent paths (\cref{sec:cond}),
an independent check of the data in \cref{sec:results}.

\section{Results}
\label{sec:results}

\noindent\textbf{Data provenance.} The tables of this section come from several
batches of runs, each drawing its own random blocks, and the archived terminal
output of every batch is distributed with the reproduction material together with
the mapping from table to output file. Comparisons should therefore be made
\emph{within} a table rather than across tables at a precision finer than a few
parts in a thousand: the same nominal configuration can differ by more than one
standard error, as for ReLU at $L_m=2$, $n=512$, where the value is $1.124$ in
\cref{tab:infl} and $1.110$ in \cref{tab:srfull}, about $0.9\%$. None of the
qualitative conclusions of this section depends on a difference of that size.

\subsection{Failure of the zero-mean assumption}
\label{sec:bias}

We test (H1) with a $z$ statistic on the block mean: $N=4\times10^4$ blocks and
$z=\bar e_{\text{block}}/\bigl(\hat\sigma_{\text{block}}/\sqrt N\bigr)$.
\Cref{fig:biasvar} (top) gives the results.

\begin{figure}[!tb]
\centering
\includegraphics[width=0.92\textwidth]{figures/fig_biasvar.pdf}
\caption{Top: significance $|z|$ of the block-mean bias (logarithmic vertical
axis); the red dashed line is the threshold $|z|=4$, blue axis labels denote
symmetric distributions and red ones asymmetric distributions. Bottom: ratio of
the measured variance to the classical prediction $\mathbb{E}[s^2]/12$; the
dashed line is the classical prediction $1$. In both panels all four mantissa
widths lie below the threshold for the symmetric distributions, and all
asymmetric distributions are significant for $L_m\le4$. The ReLU variance ratio
is constantly about $0.506$; for the spike mixture it drops to $0.093$ at low
precision (the classical formula overestimates by about a factor of $10$).}
\label{fig:biasvar}
\end{figure}

All symmetric distributions pass (H1) with $|z|\le1.63$, whereas all asymmetric
ones fail for $L_m\le4$ ($z=-126.80$, $-60.57$ and $-39.48$ at $L_m=2$ for the
lognormal, half-normal and ReLU respectively, and $-19.92$, $-15.22$, $-9.26$ at
$L_m=4$), and the bias grows as the mantissa width decreases. The relative bias
(bias divided by the mean magnitude) reaches $-6.19\%$ for the lognormal at
$L_m=2$. Values of $|z|$ for $L_m=3$ are $-61.69$, $-30.96$ and $-19.45$; all
figures are those plotted in \cref{fig:biasvar} (top), computed from
$N=4\times10^4$ blocks.

The mechanism agrees with \cref{eq:mxi}: the expected round-to-nearest error
vanishes for symmetric inputs, whereas for nonnegative inputs the opportunities to
round down and up are unequal and a systematic negative bias appears. Post-ReLU
activations are exactly nonnegative---not a special case, but the norm in the
forward pass of modern neural networks.

\subsection{Structural bias of the classical variance formula}
\label{sec:var}

Taking $\mathbb{E}[s^2]/12$ as the prediction and the measured variance as the
comparison, the ratios are given in \cref{tab:var} and \cref{fig:biasvar}
(bottom).

\begin{table}[!tb]
\centering
\caption{Ratio of the measured variance to the classical prediction
$\mathbb{E}[s^2]/12$.}
\label{tab:var}
\small
\begin{tabular}{@{}lrrrr@{}}
\toprule
Distribution & $L_m=2$ & $L_m=3$ & $L_m=4$ & $L_m=7$ \\
\midrule
Gaussian           & $0.999$ & $0.999$ & $0.999$ & $1.001$ \\
Student-$t_3$      & $0.808$ & $0.913$ & $0.914$ & $1.002$ \\
Spike mixture      & $\mathbf{0.093}$ & $\mathbf{0.148}$ & $\mathbf{0.319}$ & $1.002$ \\
Lognormal          & $0.854$ & $1.011$ & $1.043$ & $1.002$ \\
Half-normal        & $0.996$ & $0.999$ & $1.000$ & $1.001$ \\
ReLU               & $\mathbf{0.506}$ & $\mathbf{0.507}$ & $\mathbf{0.506}$ & $\mathbf{0.506}$ \\
\bottomrule
\end{tabular}
\end{table}

By \cref{prop:var}, the ratio of the measured to the classical variance has the
closed form
\begin{equation}
\frac{\operatorname{Var}(e)}{\mathbb{E}[s^2]/12}
=\frac{12\bigl(\mathbb{E}[s^2\gamma(\xi)]-(\mathbb{E}[e])^2\bigr)}{\mathbb{E}[s^2]} ,
\label{eq:ratio}
\end{equation}
so a deviation from $1$ has a single source: the departure of
$\gamma(\xi)=\mathbb{E}[g(x/s)^2\mid\xi]$ from $1/12$ (plus the small term
$-12(\mathbb{E}e)^2/\mathbb{E}[s^2]$). The condition $\gamma\equiv1/12$ holds
when, given $\xi$, the normalized residual is \emph{conditionally uniform}
within its rounding cell---assumed implicitly by the classical model, and false
for neural-network activations. Three departures appear.

First, the ReLU ratio is pinned near $0.500$ and does not vary with the mantissa
width or the block length, because the model omits a \emph{probability atom}:
exactly half of the ReLU outputs are precisely zero and their rounding error is
identically zero, so
\begin{equation}
\gamma(s)=\bigl(1-p_0\bigr)\cdot\tfrac1{12}+p_0\cdot 0=\tfrac1{24},
\qquad p_0=\mathbb{P}(x=0)=\tfrac12 ,
\label{eq:atom}
\end{equation}
which predicts a ratio $=1-p_0=0.500$ (that the non-atomic part still has an
approximately conditionally uniform residual is consistent with the half-normal
measurements). The measured value is $0.506$ (\cref{tab:var},
$4\times10^4$ blocks), and all values over $L_m\in\{2,3,4,7\}$ and
$n\in\{32,512,4096\}$ lie in $0.499$--$0.509$; $p_0=1/2$ is fixed by the
definition of ReLU and is not a fitted parameter.

Second, the same explanation predicts the ratio under stochastic rounding. By
\cref{cor:srvar}, stochastic rounding changes the within-cell variance from
$1/12$ to $1/6$, so
\begin{equation}
\frac{\operatorname{Var}(e)}{\mathbb{E}[s^2]/12}\Big|_{\text{SR}}
=2(1-p_0)=1.000\ \text{(ReLU)},\qquad =2.000\ \text{(half-normal, }p_0=0\text{)} .
\end{equation}
The data are exactly of this form: $1.015$ for ReLU and $2.002$ for the
half-normal (from the rounding-rule runs of \cref{sec:modes}), and $2.000$ for the
Gaussian in the measurement
of \cref{sec:cond} ($L_m=4$, $n=4096$). This is a cross-validation with no
parameter tuned after the fact.

Third, the value $0.093$ for the spike mixture and $0.506$ for ReLU come from the
same mechanism: $\gamma$ is depressed by the \emph{dead zone}. When the step is
much larger than the bulk scale of the data, most elements satisfy $|u|<1/2$ and
round to zero, and the residual $g(u)=-u$ is tiny; the binned data of
\cref{sec:cond} show that $96\%$ of the spike-mixture elements fall in the dead
zone at $n=4096$, where the within-bin average of $\gamma$ is only $0.0045$. In
unified form: the classical formula systematically overestimates the error
variance when the dead zone dominates, whereas when $s$ is comparable to the bulk
scale of the data, $\gamma$ can slightly exceed $1/12$ (for instance $1.043$ for
the lognormal at $L_m=4$) and the ratio slightly exceeds $1$.

The deviations of this subsection and those of \cref{sec:infl} are mutually
independent: by \cref{prop:var} the former is controlled by
$\gamma\not\equiv1/12$ and the latter by $\delta\not\equiv0$, and
\cref{sec:sr} shows that stochastic rounding turns the ratio studied here into
$2.000$ (\cref{cor:srvar}) while the inflation vanishes simultaneously, so the
two must be corrected separately. One methodological caveat: both numerator and
denominator of \cref{eq:ratio} are dominated by rare large-$\xi$ blocks, so
convergence is slow---discarding sparse $\xi$ bins with fewer than $20$ blocks
changes $\operatorname{Var}(e)$ by up to $4.9\%$, which explains the discrepancy
between $0.854$ in \cref{tab:var} and $0.880$ in \cref{sec:model} for the same
nominal configuration (lognormal, $L_m=2$, $n=32$).

\subsection{Scaling behaviour of the variance inflation}
\label{sec:infl}

The variance inflation factor is the central quantity here. We estimate
$\mathcal I_n$ by \cref{eq:estimator}, with the block length $n$ increasing from
$32$ to $4096$. \Cref{fig:infl} and \cref{tab:infl} give the results.

\begin{figure}[!tb]
\centering
\includegraphics[width=0.88\textwidth]{figures/fig3_inflation.pdf}
\caption{Variance inflation factor $\mathcal I_n$ as a function of the block
length $n$ (log--log axes). Left: $L_m=2$; right: $L_m=4$. The grey dotted line
is $\propto n$ and the grey dash-dotted line is $\propto\sqrt n$. The symmetric
control (Gaussian, dashed) is identically $1$ for all $n\le4096$. The asymmetric
distributions grow systematically with $n$, and the lognormal reaches a tail
scaling exponent of $1.134$ at $L_m=4$.}
\label{fig:infl}
\end{figure}

\begin{table}[!tb]
\centering
\caption{Variance inflation factor $\mathcal I_n$ ($20$ groups $\times$ $1000$
blocks). The last column is the local log-slope over the largest measured
interval ($n\in[512,4096]$ for $L_m=2$, $n\in[2048,4096]$ for $L_m=4$).}
\label{tab:infl}
\setlength{\tabcolsep}{4pt}
\footnotesize
\begin{tabular}{@{}llrrrrrrc@{}}
\toprule
$L_m$ & Distribution & $32$ & $128$ & $512$ & $1024$ & $2048$ & $4096$ & Tail slope \\
\midrule
$2$ & Gaussian (control) & $1.000$ & $0.985$ & $0.997$ & $1.009$ & $1.026$ & $1.013$ & $\approx0$ \\
$2$ & Lognormal       & $3.075$ & $9.809$ & $31.341$ & $51.480$ & $102.555$ & $\mathbf{176.996}$ & $0.833$ \\
$2$ & ReLU            & $1.037$ & $1.007$ & $1.124$ & $1.520$ & $2.679$ & $\mathbf{6.132}$ & $0.816$ \\
$2$ & Half-normal     & $1.015$ & $1.043$ & $1.441$ & $2.657$ & $6.321$ & $\mathbf{14.989}$ & $1.126$ \\
\midrule
$4$ & Gaussian (control) & $0.996$ & $0.984$ & $1.009$ & $0.987$ & $1.013$ & $0.989$ & $\approx0$ \\
$4$ & Lognormal       & $1.328$ & $3.511$ & $15.601$ & $34.748$ & $76.063$ & $\mathbf{165.025}$ & $1.134$ \\
$4$ & ReLU            & $1.005$ & $0.986$ & $1.026$ & $1.030$ & $1.084$ & $1.300$ & $0.114$ \\
$4$ & Half-normal     & $0.996$ & $1.020$ & $1.013$ & $1.096$ & $1.284$ & $1.834$ & $0.286$ \\
\bottomrule
\end{tabular}
\end{table}

The measurements support three conclusions.

\paragraph{(i) (H1) and (H3) hold for symmetric distributions} For the Gaussian,
$\mathcal I_n=1.000\pm0.01$ and stays flat in $n$ up to $n=4096$; the same holds at
every measured block length for Student-$t_3$ and the spike mixture, whose
$\mathcal I_n$ lies in $0.99$--$1.01$ throughout (\cref{app:full}).
This explains why the hardware literature did not go wrong in its original,
approximately symmetric setting, and it rules out implementation artefacts. It
does not mean that the classical model holds as a whole: \cref{sec:var} shows that
the symmetric Student-$t_3$ has a variance ratio as low as $0.808$ for
$L_m\le4$, so (H2) can still fail. A value $\mathcal I_n=1$ only excludes
correlation; it does not guarantee the variance formula.

\paragraph{(ii) Growth on asymmetric distributions is unsaturated, but the
exponent cannot exceed $1$} For the lognormal at $L_m=4$ the local log-slope over
$n\in[512,4096]$ is $1.134$, but by \cref{cor:bounds} $\mathcal I_n\le n$, so the
\emph{asymptotic} slope is at most $1$: the value $1.134$ can only be a transient,
and extrapolating $n^{1.13}$ would contradict the rigorous bound
(\cref{sec:open} proves $\mathcal I_n/n\to0$). The reduced model \cref{eq:reduced2}
gives the right reading: \cref{tab:infl} yields $\rho_n=0.043$ at $n=4096$ for the
lognormal, $0.0013$ for ReLU and $0.0034$ for the half-normal, so over the measured
range the evidence supports only ``unsaturated''. By \cref{sec:domains-res} the
lognormal value of $\rho_n$ is already decaying at $n=65536$, so we claim nothing
about the asymptotic exponent here; \cref{sec:inflbig} and \cref{sec:open} give the
complete conclusion.

\paragraph{(iii) The ``threshold-like'' behaviour of ReLU is a misjudgement} At
$L_m=2$ the value of $\mathcal I_n$ for ReLU is close to $1$ for $n\le128$ and
rises afterwards, so a finite range alone suggests a ``threshold effect''; but the
local log-slope increases monotonically with $n$:
$-0.030,\ -0.014,\ 0.030,\ 0.130,\ 0.435,\ 0.818,\ \mathbf{1.195}$, which is a
transition from below the detection limit to superlinear growth, not a threshold
tending to zero.

\paragraph{Non-monotonicity of the mantissa-width dependence} The lognormal has a
tail slope of $0.833$ at $L_m=2$ and $1.134$ at $L_m=4$, whereas ReLU and the
half-normal behave in the opposite way, with almost no effect at $L_m=4$ and a
marked one at $L_m=2$. This follows from the dead-zone law
$\rho\propto4^{-L_m}$ (\cref{eq:lmscaling}): each extra mantissa bit shrinks
$\mathcal I_n-1$ by a factor of $4$, and the data give ratios of $16.3$ (ReLU) and
$16.9$ (half-normal), consistent with $4^2=16$ for two mantissa bits. The
lognormal value of $\rho$ is nearly the same at the two widths (ratio $1.12$)
because its step lies near the peak of $|\delta(s)|$, where the scaling law does
not apply; a rigorous characterization of $\delta$ in that transition region
remains open.

\subsection{Sensitivity to the rounding rule}
\label{sec:modes}

\begin{table}[!tb]
\centering
\caption{Sensitivity to the rounding rule ($n=32$, $L_m=4$; the stochastic
rounding column uses a single group of $1000$ blocks, whose error bars are far
larger than the $20\times1000$ blocks used for the definitive conclusions, so its
$\mathcal I_{32}$ serves only as a qualitative reference; for the definitive test
see \cref{tab:srfull} in \cref{sec:sr}).}
\label{tab:modes}
\setlength{\tabcolsep}{4pt}
\small
\begin{tabular}{@{}llrrr@{}}
\toprule
Distribution & Rounding rule & $z_{\text{bias}}$ & Measured/predicted variance & $\mathcal I_{32}$ \\
\midrule
\multirow{4}{*}{Lognormal}
 & nearest (RNE) & $-25.33$ & $1.043$ & $1.253$ \\
 & half-away & $-24.68$ & $1.038$ & $1.188$ \\
 & truncation & $\mathbf{-422.71}$ & $1.623$ & $14.270$ \\
 & \textbf{stochastic} & $\mathbf{-0.98}$ & $2.066$ & $0.825$ \\
\midrule
\multirow{4}{*}{ReLU}
 & nearest (RNE) & $-11.29$ & $0.508$ & $1.196$ \\
 & half-away & $-13.16$ & $0.507$ & $1.072$ \\
 & truncation & $\mathbf{-603.77}$ & $1.315$ & $2.931$ \\
 & \textbf{stochastic} & $\mathbf{-0.81}$ & $1.015$ & $1.079$ \\
\midrule
\multirow{4}{*}{Half-normal}
 & nearest (RNE) & $-21.08$ & $0.999$ & $1.109$ \\
 & half-away & $-18.23$ & $1.001$ & $1.182$ \\
 & truncation & $\mathbf{-982.14}$ & $1.155$ & $5.357$ \\
 & \textbf{stochastic} & $\mathbf{0.97}$ & $2.002$ & $1.066$ \\
\bottomrule
\end{tabular}
\end{table}

To rule out that the conclusions are merely an artefact of round-to-nearest, we
compared four rounding rules at $n=32$ and $L_m=4$: round-to-nearest,
round-half-away-from-zero, truncation toward zero, and stochastic rounding
(\cref{tab:modes}). Two conclusions follow, and both are visible in
\cref{fig:srmodes} (right). First, the finding of \cref{sec:bias} is not an
artefact of round-to-nearest: round-half-away gives almost the same bias,
$-25.33$ versus $-24.68$ for the lognormal (and $-11.29$ versus $-13.16$ for
ReLU). Second, stochastic rounding removes the bias completely
($-25.33\to-0.98$ for the lognormal, $-11.29\to-0.81$ for ReLU, $-21.08\to0.97$
for the half-normal), consistent with \cref{prop:sr}, while truncation amplifies
it sharply ($-422.71$, $-603.77$ and $-982.14$ respectively), in agreement with
the engineering consensus that truncation is unusable. The same runs give a
self-check on the variance ratio: under SR it is $2.000$ for the Gaussian and
$2.002$ for the half-normal, the value corresponding to the theoretical SR error
variance $s^2/6=2\cdot s^2/12$, which shows that the $s^2/12$ of the classical
formula is specific to round-to-nearest. Their $\mathcal I_{32}$ values under SR
($0.825$--$1.079$) deviate noticeably from $1$ because a single group of $1000$
blocks is used at $n=32$; the definitive test of \cref{prop:sr} uses
$20\times1000$ blocks and gives $24/24$ configurations within
$1.000\pm2.3\,\mathrm{se}$ for $n\ge512$ (\cref{tab:srfull}).

\subsection{A falsifiable test of the mechanism}
\label{sec:sr}

\begin{figure}[!tb]
\centering
\includegraphics[width=0.92\textwidth]{figures/fig_srmodes.pdf}
\caption{A falsifiable test of the mechanism, and the rounding-rule control.
Left: comparison of $\mathcal I_n$ under round-to-nearest and stochastic rounding
at $n=4096$, $L_m=2$ (logarithmic vertical axis; the black dashed line is the
predicted value $1$). The same distributions, the same block lengths and the
same code path, with only the rounding rule changed: the lognormal drops from
$176.05$ to $0.988$. Right: significance $|z|$ of the block-mean bias under four
rounding rules ($n=32$, $L_m=4$, logarithmic axis). Round-half-away gives almost
the same bias as round-to-nearest, truncation amplifies it sharply, and
stochastic rounding removes it completely.}
\label{fig:srmodes}
\end{figure}

\Cref{prop:sr} yields a high-risk inference about the mechanism: after switching
to stochastic rounding, $\mathcal I_n$ must equal $1$ exactly, independently of
the distribution, the block length and the mantissa width; any significant
deviation would refute \cref{prop:decomp} or \cref{prop:decay}.

\begin{table}[!tb]
\centering
\caption{Test of the mechanism at $n=4096$: is $\mathcal I_n$ identically $1$
under stochastic rounding (the predicted value)? Standard errors are in
parentheses; $\Delta$ is the deviation of SR from $1$ in standard errors. The
same comparison at $n=32$ and $n=512$ is reported in the text.}
\label{tab:srfull}
\setlength{\tabcolsep}{4pt}
\footnotesize
\begin{tabular}{@{}llrrr@{}}
\toprule
$L_m$ & Distribution & RNE & SR & $\Delta$ \\
\midrule
$2$ & Gaussian    & $1.004\,(0.009)$ & $0.988\,(0.012)$ & $-1.0$ \\
$2$ & Lognormal   & $\mathbf{176.051}\ (1.373)$ & $\mathbf{0.988}\,(0.012)$ & $-1.0$ \\
$2$ & ReLU        & $\mathbf{6.111}\,(0.074)$ & $\mathbf{1.009}\,(0.012)$ & $0.8$ \\
$2$ & Half-normal & $\mathbf{14.871}\,(0.094)$ & $\mathbf{0.994}\,(0.012)$ & $-0.5$ \\
\midrule
$4$ & Gaussian    & $1.009\,(0.010)$ & $0.975\,(0.011)$ & $-2.3$ \\
$4$ & Lognormal   & $\mathbf{168.274}\ (3.073)$ & $\mathbf{0.991}\,(0.014)$ & $-0.6$ \\
$4$ & ReLU        & $1.290\,(0.014)$ & $1.000\,(0.010)$ & $0.0$ \\
$4$ & Half-normal & $1.775\,(0.023)$ & $0.991\,(0.013)$ & $-0.7$ \\
\bottomrule
\end{tabular}
\end{table}

\noindent\textbf{Result.} All $24$ stochastic rounding configurations fall within
$1.000\pm3\,\mathrm{se}$ ($0.974$--$1.061$). \Cref{tab:srfull} gives the values at
$n=4096$; at $n=32$ and $n=512$ the SR values are, at $L_m=2$, $1.013$, $0.983$,
$1.004$, $1.013$ and $1.003$, $1.013$, $0.996$, $1.000$ (Gaussian, lognormal,
ReLU, half-normal), and at $L_m=4$, $0.992$, $1.061$, $1.005$, $1.004$ and
$1.010$, $1.036$, $1.003$, $0.974$. The corresponding RNE values are $1.010$,
$3.134$, $1.012$, $1.007$ and $0.999$, $31.507$, $1.110$, $1.473$ at $L_m=2$,
and $1.015$, $1.342$, $1.012$, $1.023$ and $1.009$, $15.798$, $0.996$, $1.039$ at
$L_m=4$. The contrast is unmistakable: changing only the rounding rule drops
$\mathcal I_n$ at $n=4096$ from $176$ to $0.99$.

\subsection{Direct decomposition conditioned on $\xi$, and a test of
assumption (A2)}
\label{sec:cond}

The identity of \cref{prop:decomp} splits the covariance into two terms, and
(A2) asserts that the second is everything. We measure both directly, as in
\cref{sec:infl} ($20$ groups $\times$ $1000$ blocks): bins are the integer values
of $\xi$, centred by the \emph{within-bin mean} $m_b$ rather than the within-block
mean (see \cref{sec:estimator}), keeping only bins with at least $20$ blocks
(retained mass $\ge0.999$).

\begin{table}[!tb]
\centering
\caption{Direct measurement of the conditional decomposition ($L_m=2$,
$n=4096$, round-to-nearest). $\operatorname{Cov}$ is the global pairwise
covariance; $\operatorname{Var}(m)$ and $\mathbb{E}[c]$ are the two terms
estimated by binning on $\xi$; the last column is
$\mathbb{E}[c(\xi)]/\operatorname{Cov}$.}
\label{tab:cond}
\footnotesize
\begin{tabular}{@{}lrrrrr@{}}
\toprule
Distribution & $\mathcal I_{4096}$ & $\operatorname{Cov}(e_i,e_j)$ & $\operatorname{Var}(m(\xi))$ & $\mathbb{E}[c(\xi)]$ & $\mathbb{E}[c]/\operatorname{Cov}$ \\
\midrule
Gaussian      & $0.988$  & $-1.02\times10^{-7}$ & $-7.3\times10^{-10}$ & $-1.02\times10^{-7}$ & $0.99$ \\
Student-$t_3$ & $0.997$  & $-9.23\times10^{-7}$ & $-4.5\times10^{-8}$  & $-8.78\times10^{-7}$ & $0.95$ \\
Spike mixture & $0.985$  & $-4.30\times10^{-9}$ & $-6.5\times10^{-11}$ & $-4.24\times10^{-9}$ & $0.99$ \\
\midrule
Lognormal     & $\mathbf{175.77}$ & $6.79\times10^{-2}$ & $6.79\times10^{-2}$ & $5.5\times10^{-9}$ & $8.1\times10^{-8}$ \\
ReLU          & $\mathbf{6.163}$  & $1.80\times10^{-5}$ & $1.79\times10^{-5}$ & $4.5\times10^{-8}$ & $2.5\times10^{-3}$ \\
Half-normal   & $\mathbf{14.880}$ & $1.19\times10^{-4}$ & $1.19\times10^{-4}$ & $4.6\times10^{-8}$ & $3.9\times10^{-4}$ \\
\bottomrule
\end{tabular}
\end{table}

\Cref{tab:cond} reports the measurements; four findings follow.

\paragraph{(1) Three paths agree; two identities at machine precision} For each
configuration $\mathcal I_n$ is produced in three ways: by \cref{eq:estimator},
by $1+(n-1)\rho$ recovery from the global pairwise covariance, and by
$1+(n-1)\operatorname{Var}(m)/\operatorname{Var}(e)$ from binning. The first two
agree to four decimal places over all $72$ configurations; over the $11$ with
inflation ($\mathcal I_n>1.2$) the binned and direct estimates differ by at most
$5.7\%$ (lognormal, $L_m=4$, $n=512$), the remainder within $3\%$. On the retained
bins the relative closure error of
$\operatorname{Var}(e)=\mathbb{E}[\operatorname{Var}(e\mid\xi)]+\operatorname{Var}(m(\xi))$
is $3.4\times10^{-16}$, and that of the block-mean decomposition
$\operatorname{Var}(\bar e)=\operatorname{Var}(m)+\mathbb{E}[\operatorname{Var}(e\mid\xi)]/n+(n-1)\mathbb{E}[c(\xi)]/n$
is $2.2\times10^{-15}$; its third term is exactly the contribution of the
violation of (A2) and also decomposes the $z$ statistic of \cref{sec:bias}.

\paragraph{(2) Share from the violation of (A2); numerical validity of (A2)} The
share $(n-1)\mathbb{E}[c(\xi)]/\operatorname{Var}(e)$ of $\mathcal I_n-1$ from
$\mathbb{E}[c(\xi)]$ is at most $1.6$ over all $36$ round-to-nearest
configurations, while the inflated ones reach $\mathcal I_n-1=175$: it is $1.1\%$
in the worst of the $11$ configurations with $\mathcal I_n>1.2$ and never exceeds
$2\%$ under substantial inflation. The ratio
$\mathbb{E}[c(\xi)]/\operatorname{Cov}$ is $8\times10^{-8}$ (lognormal),
$2.5\times10^{-3}$ (ReLU) and $3.9\times10^{-4}$ (half-normal); for the three
symmetric distributions it is meaningless, $\operatorname{Cov}$ already lying below
$10^{-6}$. Essentially all inflation therefore comes from
$\operatorname{Var}(m(\xi))$, which is the numerical content of (A2).

\paragraph{(3) The upper bound holds; a candidate lower bound is refuted} All $72$
configurations satisfy $0\le\mathcal I_n\le n$ (\cref{cor:bounds}). The candidate
lower bound $\mathcal I_n\ge n(\mathbb{E}[e])^2/\operatorname{Var}(e)$ \emph{does
not hold}: for the lognormal at $L_m=2$, $n=4096$ its right-hand side is $1586$
while the measured $\mathcal I_n=175.8$, and $6$ of the $36$ round-to-nearest
configurations violate it. The bias $\mathbb{E}[e]=\mathbb{E}[m(\xi)]$ and the
within-block correlation
$\operatorname{Cov}=\operatorname{Var}(m(\xi))+\mathbb{E}[c(\xi)]$ are both moments
of the same $m(\xi)$, testing only $m\equiv0$; neither being nonzero implies the
other is.

\paragraph{(4) Stochastic rounding control} Under stochastic rounding both terms
fall to numerical zero ($|\operatorname{Var}(m)|\le4\times10^{-6}$,
$|\mathbb{E}[c]|\le2.6\times10^{-4}$, every configuration with $\mathcal I_n$
within $1.000\pm0.03$), consistent with \cref{prop:sr}: $\mathbb{E}[c(\xi)]=0$ is
not an extra assumption, holding automatically when $m\equiv0$.

\subsection{Parameter-free predictions of the reduced model}
\label{sec:model}

Finally we test the reduced model \cref{eq:reduced,eq:reduced2} of
\cref{sec:modeltheory}. Per configuration we estimate the one-dimensional
$\delta(s),\gamma(s)$ from \cref{eq:deltagamma} on $4\times10^6$ samples
(independently of any blocking), take the empirical distribution of $s$ (the
order statistic of the block maximum) and compute $\mathbb{E}[e]$,
$\operatorname{Var}(e)$ and $\mathcal I_n$ from \cref{eq:reduced2}; apart from
the sample size there is no adjustable parameter, and $\delta,\gamma$ are not
fitted to the simulation data.

\begin{table}[!tb]
\centering
\caption{Parameter-free predictions of the reduced model against measurement
(round-to-nearest, $20$ groups $\times$ $1000$ blocks). Except for the two
$n=32$ entries of the lognormal, which are given because they are the largest
errors at small $n$, the table lists the $n=4096$ configurations; the model also
predicts $\mathbb{E}[e]$ and $\operatorname{Var}(e)$ (not listed; maximum
relative error $<2\%$).}
\label{tab:model}
\footnotesize
\begin{tabular}{@{}llrrr@{}}
\toprule
$L_m$ & Distribution & $\mathcal I_n$ measured & $\mathcal I_n$ model & Relative error \\
\midrule
$2$ & Gaussian    & $0.9883$ & $1.0003$ & $+1.2\%$ \\
$2$ & Student-$t_3$ & $0.9965$ & $1.0005$ & $+0.4\%$ \\
$2$ & Lognormal, $n=32$ & $3.1354$ & $3.2013$ & $+2.1\%$ \\
$2$ & Lognormal, $n=4096$ & $175.7679$ & $176.4593$ & $+0.4\%$ \\
$2$ & Spike mix.  & $0.9853$ & $1.0000$ & $+1.5\%$ \\
$2$ & ReLU        & $\mathbf{6.1631}$ & $6.3155$ & $+2.5\%$ \\
$2$ & Half-normal & $\mathbf{14.8804}$ & $15.0719$ & $+1.3\%$ \\
\midrule
$4$ & Gaussian    & $1.0014$ & $1.0019$ & $+0.1\%$ \\
$4$ & Student-$t_3$ & $0.9888$ & $1.0010$ & $+1.2\%$ \\
$4$ & Lognormal   & $157.1972$ & $156.7585$ & $-0.3\%$ \\
$4$ & Spike mix.  & $0.9959$ & $1.0000$ & $+0.4\%$ \\
$4$ & ReLU        & $\mathbf{1.3239}$ & $1.3646$ & $+3.1\%$ \\
$4$ & Half-normal & $\mathbf{1.8243}$ & $1.8868$ & $+3.4\%$ \\
\bottomrule
\end{tabular}
\end{table}

\noindent Over all $36$ round-to-nearest configurations the relative deviation
$\mathcal I_n^{\text{model}}/\mathcal I_n^{\text{measured}}-1$ is at most $4.47\%$
in absolute value (median $1.15\%$), with $\mathbb{E}[e]$ and
$\operatorname{Var}(e)$ predicted to below $2\%$, mostly $0.5\%$; the $n=32$
configurations, at which the model is at its least accurate, are within the same
bound. The two-orders-of-magnitude growth in \cref{fig:infl} is thus a direct
output of \cref{eq:reduced2}, not an empirical regularity requiring a separate
fit.

\medskip
\noindent\textbf{A boundary that must be observed.} The two readings of $\rho_n$
--- the decomposition $\rho_n=\operatorname{Var}(m(\xi))/\operatorname{Var}(e)$
and the ratio $\mathbb E[s^2\delta(s)^2]/\mathbb E[s^2\gamma(s)]$ of direct
expectations over $s$ --- differ by the same subtraction of $(\mathbb E e)^2$ in
numerator and denominator. That is negligible for most configurations here
($(\mathbb E e)^2/\operatorname{Var}(e)\approx5\times10^{-6}$ for the lognormal at
$L_m=2$, $n=4096$), but for Pareto$(4)$, where
$\mathbb E[m]\approx-0.44$ is not small, the two are not interchangeable and one
must use the decomposition, as the ``model'' values here and in
\cref{tab:bounded} do.

\medskip
\noindent\textbf{Two structural facts.} First, $\delta$ and $\gamma$ in
\cref{eq:deltagamma} depend only on $s$, not on $L_m$: the lognormal gives the
same $s=4$ at $(L_m,\xi)=(2,4)$ and $(4,6)$, with binned
$\delta=-0.11221$ and $-0.11223$, and at $(2,5)$ and $(4,7)$, both $s=8$, the
values $-0.11461$ and $-0.11462$, supporting the main-regime conclusion of
\cref{lem:cond}. Second, the dead-zone law is linear: with our sign convention
$\delta=\mathbb{E}[\mathrm{RN}(x/s)-x/s]$, hence negative,
$\delta(0.125)=-0.00414$, $\delta(0.25)=-0.00841$, $\delta(0.5)=-0.01668$ and
$\delta(1)=-0.03434$, that is $\delta(s)=-0.0334\,s$. For ReLU, $\delta$ is
exactly half the half-normal value ($-0.0167\,s$: a truncated half-normal minus
an atom of $p_0=1/2$, the standard half-normal having
$\delta=-0.0334\,s$), while $12\gamma\equiv1.000$ (half-normal) and
$0.500$ (ReLU) in that range. Hence $\rho\propto4^{-L_m}$ in
\cref{eq:lmscaling}.

\subsection{Pushing the block length further}
\label{sec:inflbig}

Beyond $n\approx10^5$ for Pareto$(4)$, the measured $\mathcal I_n$ no longer
tracks a constant-$\rho_n$ extrapolation. \Cref{tab:bign} records the large-$n$
measurements and the corresponding population prediction of the reduced model.

\begin{table}[!tb]
\centering
\caption{Measured $\mathcal I_n$ for Pareto$(4)$, $L_m=2$ (independently
re-implemented). Rows 1--3 vary $n$ with $4000/3000/30000/4000$ blocks; row 4
holds $n=262144$ fixed while the number of blocks grows through
$10^3,3\times10^3,9\times10^3,2.7\times10^4$. The last row gives two independent
measurements at $n=10^6$ with $1200$ and $30000$ blocks.}
\label{tab:bign}
\small
\begin{tabular}{@{}lrrrr@{}}
\toprule
$n$ & $64$k & $256$k & $1$M & $10$M \\
\midrule
$\mathcal I_n$ & $45426$ & $41421$ & $36065$ & $9622$ \\
$\mathcal I_n/\sqrt n$ & $177.4$ & $80.9$ & $36.1$ & $3.04$ \\
$\mathcal I_n/n$ ($n$ fixed) & $0.125$ & $0.158$ & $0.146$ & $0.135$ \\
\midrule
$n=1$M, $1200$ blocks & \multicolumn{4}{c}{$\mathcal I_n=37692$ ($4.5\%$ from $36065$ with $30000$ blocks)} \\
\bottomrule
\end{tabular}
\end{table}

As $n$ increases from $64$k to $10^7$ (a factor of $1.6\times10^5$), $\mathcal I_n$
decreases from $4.5\times10^4$ to $9.6\times10^3$ and $\mathcal I_n/\sqrt n$ from
$177$ to $3.0$ (so $\mathcal I_n/n$ falls from $0.693$ to $0.00096$). A constant
$\rho_n\approx0.70$ would predict multiplication by $16$ when $n$ grows by $16$,
and that behaviour holds only for $n\le65536$; comparing $n=10^6$ with $n=10^7$
gives a local log-slope $\log(36065/9622)/\log 10=0.57$, consistent with
$\mathcal I_n\le n$ (\cref{cor:bounds}) and ruling out asymptotically linear
growth. \Cref{prop:decay} states that $B_n=\operatorname{Var}(m(\xi))\to0$ and
$\rho_n\to0$, hence $\mathcal I_n=o(n)$: linear accumulation is impossible for any
unbounded distribution with finite second moment. Evaluating \cref{eq:reduced2}
exactly by one-dimensional quadrature (see the remark in \cref{sec:open}) gives
$3.61\times10^4$ at $n=10^6$ and $9.2\times10^3$ at $n=10^7$, matching the measured
$3.61\times10^4$ and $9.6\times10^3$: the fall-off is the decline of the population
curve past its peak, not an artefact of insufficient sampling, and the row with $n$
fixed shows stability, $\mathcal I_n/n$ staying within $0.125$--$0.158$ as the
block count grows by a factor of $27$. This paper claims only that the independence
premise fails for asymmetric inputs, and claims no particular asymptotic order of
accumulation: what is a theorem is $\mathcal I_n\le n$ (\cref{cor:bounds}),
$\mathcal I_n\to1$ for bounded support (\cref{prop:bound}) and $\mathcal I_n=o(n)$
for unbounded support (\cref{prop:decay}).

\subsection{Numerical comparison of bounded and unbounded support}
\label{sec:domains-res}

This subsection separates the two factors ``bounded versus unbounded'' and
``light versus heavy tails'' using five new distributions: Beta$(2,5)$,
Beta$(5,2)$ and a half-normal truncated to $[0,2]$ (bounded, asymmetric), the
exponential distribution (unbounded, light-tailed) and Pareto$(4)$ (unbounded,
regularly varying tail). The methodology is that of \cref{sec:cond}.

\begin{table}[!tb]
\centering
\caption{Bounded versus unbounded support ($L_m=2$). The bounded
distributions are supported on $[0,1]$, $[0,1]$ and $[0,2]$ respectively; all
others are unbounded. Values of $\mathbb{P}(\xi=\text{mode})$ and of the model
$\mathcal I_n$ are taken at $n=4096$.}
\label{tab:bounded}
\footnotesize
\setlength{\tabcolsep}{4pt}
\begin{tabular}{@{}lrrrrr@{}}
\toprule
Distribution & $n=32$ & $n=512$ & $n=4096$ & $\mathbb{P}(\xi{=}\text{mode})$ & $\mathcal I_n$ model \\
\midrule
Beta$(2,5)$      & $1.001$ & $1.007$ & $1.009$ & $1.000$ & $1.000$ \\
Beta$(5,2)$      & $1.014$ & $1.007$ & $0.999$ & $1.000$ & $1.000$ \\
Truncated half-normal & $0.999$ & $0.993$ & $0.995$ & $1.000$ & $1.000$ \\
\midrule
Exponential      & $1.235$ & $8.284$ & $41.885$ & $0.747$ & $41.448$ \\
Pareto$(4)$      & $7.154$ & $290.22$ & $2886.30$ & $0.573$ & $2887.29$ \\
Lognormal        & $3.197$ & $30.537$ & $175.19$ & $0.602$ & $176.04$ \\
\bottomrule
\end{tabular}
\end{table}

\Cref{tab:bounded} fully supports \cref{prop:bound}: the three bounded-support
distributions give $\mathcal I_n=1.00\pm0.01$ for $n\le4096$, indistinguishable
from symmetric distributions and independently of the endpoint tail index. The
unbounded side is the opposite: the exponential distribution already reaches
$41.9$ at $n=4096$ and Pareto$(4)$ reaches $2886$, and the prediction error of the
reduced model is below $2\%$ for all six distributions, both widths and all three
block lengths. By \cref{prop:decay} the inflation on the unbounded side must
eventually fall back: the peak of Pareto$(4)$, $5.9\times10^4$, lies at
$n\approx10^5$ (see the remark in \cref{sec:open}), and the measurements of
\cref{sec:inflbig} agree with that curve.

\begin{table}[!tb]
\centering
\caption{$\rho_n=(\mathcal I_n-1)/(n-1)$ at large block lengths ($L_m=2$,
$20$ groups $\times$ $500$ blocks $=10^4$ blocks per configuration, averaged over
two independent seeds; the seed-to-seed difference is below $1\%$ for Pareto$(4)$
and the lognormal at all three block lengths, and reaches $5.1\%$ for the
half-normal at $n=65536$, where $\mathcal I_n$ is small and the lattice
oscillations of \cref{sec:inflbig} dominate).}
\label{tab:largen}
\footnotesize
\begin{tabular}{@{}lrrrr@{}}
\toprule
Distribution & $\rho_{4096}$ & $\rho_{16384}$ & $\rho_{65536}$ & $\mathcal I_{65536}/\mathcal I_{4096}$ \\
\midrule
Pareto$(4)$ & $0.7038$ & $0.7423$ & $0.6958$ & $\mathbf{15.77}$ \\
Lognormal   & $0.0427$ & $0.0181$ & $0.0177$ & $6.56$ \\
Half-normal & $0.0034$ & $0.0025$ & $0.0001$ & $0.60$ \\
\bottomrule
\end{tabular}
\end{table}

\Cref{tab:largen} shows two things. First, for $n\le65536$ the value of $\rho_n$
for Pareto$(4)$ is roughly constant ($0.696$--$0.742$), and a $16$-fold increase
in $n$ multiplies $\mathcal I_n$ by $15.77$, close to the factor $16.00$ predicted
by a constant $\rho_n$, i.e.
$\operatorname{Var}[\sum_ie_i]\approx0.70\,n^2\operatorname{Var}(e)$. This is the
basis of the constant-$\rho_n$ extrapolation that ``$\sqrt n$ degenerates to $n$''
over this range; \cref{sec:inflbig} shows that it is merely a \emph{transient}
over attainable block lengths. Second, the lognormal value of $\rho_n$ falls from
$0.043$ to $0.018$, making $\mathcal I_n$ markedly \emph{sublinear} (its
$\mathcal I_n/n$ at $n=65536$ is $0.0177$, against $0.0361$ at $n=10^6$ in
\cref{tab:bign}), while the half-normal exhibits an \emph{absolute} decrease of
$\mathcal I_n$ at $n=65536$ ($42.4\to9.0$) caused by the phase of $\log_2M_n$
relative to the integer grid, a jump the reduced model also reproduces. Under
unbounded support $\mathcal I_n$ therefore carries oscillations caused by lattice
alignment, and treating it as a smooth power law is dangerous.

\section{Discussion}
\label{sec:discussion}

\Cref{prop:decomp}, \cref{prop:var} and the reduced model \cref{eq:reduced} give
one picture: the within-block errors are coupled through the shared step $s$, the
rounded logarithmic scale of the largest magnitude in the block, so every
fluctuation of $s$ enters \emph{every} error with the same sign, with coefficient
the deterministic response curve $\delta(s)$. If $\delta\not\equiv0$, that common
fluctuation accumulates coherently in a summation instead of averaging away; and
since $\delta,\gamma$ are one-dimensional, all the randomness resides in the
distribution of $s$. \Cref{tab:explain} collects the mechanism's explanatory power
for the observations of \cref{sec:results}; \cref{sec:limitations} states the
boundary of validity.

\begin{table}[!tb]
\centering
\caption{Explanatory power of the mechanism for the individual observations.}
\label{tab:explain}
\setlength{\tabcolsep}{4pt}
\footnotesize
\begin{tabular}{@{}p{0.32\textwidth}p{0.44\textwidth}p{0.17\textwidth}@{}}
\toprule
Observation & Mechanistic explanation & Status \\
\midrule
Symmetric distributions give exactly $\mathcal I_n=1.0$ for $n\le4096$ & $F$ symmetric $\Rightarrow\delta(s)\equiv0$, independently of (A2) & verified (exact) \\
The within-block correlation comes from $\operatorname{Var}(m(\xi))$, not from $\mathbb{E}[c(\xi)]$; the asymptotic exponent does not exceed $1$ & exact decomposition of \cref{prop:decomp}, with the violation of (A2) accounting for at most $2\%$ of $\mathcal I_n-1$; and \cref{cor:bounds} (Cauchy--Schwarz) & verified (direct measurement); proved \\
The ``threshold-like onset'' of ReLU, and its variance ratio $0.506$ & dead-zone law $\rho\propto4^{-L_m}$ (\cref{eq:lmscaling}), measured ratio $16.3$ against $4^2=16$; atom $p_0=1/2$ gives $\gamma=1/24$ and ratio $1-p_0$, and $2(1-p_0)=1.000$ under SR & explained quantitatively \\
All numerical values of $\mathcal I_n$ & \cref{eq:reduced2}, maximum error $4.5\%$ over $36$ configurations (no fitted parameter) & verified \\
$\mathcal I_n\equiv1$ under stochastic rounding & \cref{prop:sr} ($m\equiv0$ and $c\equiv0$) & verified ($24/24$) \\
Bounded support reaches a plateau; unbounded support rises and then falls & bounded: $\xi$ concentrates on one grid point $\Rightarrow B_n=O(p^n)$; unbounded: $m(\xi)\to-\mathbb E[x_1]$ $\Rightarrow B_n\to0$ (\cref{prop:bound,prop:decay}) & proved; the peak location is given by the model \\
\bottomrule
\end{tabular}
\end{table}

\subsection{Scope of validity and practical implications}

The effect varies in strength across distributions: at the MX block length $n=32$
the lognormal gives $\mathcal I_{32}=3.08$, a summation error $1.77$ times the
independent-model prediction, whereas ReLU outputs give only $1.037$ (about
$1.8\%$). This paper therefore cannot be read as claiming that current hardware is
already seriously inaccurate. By \cref{sec:domains-res} the $\sqrt n$ advantage is
markedly weakened only for unbounded support with a \emph{regularly varying tail}
(Pareto type) and only for $n\lesssim10^5$; bounded support is better behaved than
the classical model, and fast-decaying unbounded tails such as the lognormal lie
between the two. Since \cref{prop:decay} forces the unbounded-side inflation to
fall back for large $n$, the value of the classification is in showing that block
length and tail type jointly determine the risk, not in giving an extrapolatable
power law. Finally, \cref{prop:sr} concerns the accumulation statistics of the
within-block error only: for asymmetric activations stochastic rounding removes the
bias and \emph{exactly} restores $\sqrt n$ accumulation (\cref{sec:sr}, $24/24$
configurations). This is separate from end-to-end training stability: Cim et
al.\ \cite{cim2026mxfp4} find that on native FP4 hardware, once the weight gradient
uses MXFP4, \emph{neither} stochastic rounding nor random Hadamard rotations
stabilize training, whereas a deterministic Hadamard transform does, the
instability arising from ``structured micro-scaling errors along sensitive gradient
paths, rather than insufficient stochasticity''.

\subsection{Limitations}
\label{sec:limitations}

\begin{enumerate}
\item The rate of the asymptotic decay is undetermined: rigorous are only
$\mathcal I_n\le n$ (\cref{cor:bounds}), $\mathcal I_n\to1$ for bounded support
(\cref{prop:bound}) and $\mathcal I_n=o(n)$ for unbounded support
(\cref{prop:decay}). \Cref{prop:decay} gives no rate, so $\mathcal I_n\to1$ rests on
the numerical model (see \cref{sec:open}) and the peak location is not a theorem; a
rate theorem would need a regularity assumption on $\delta(s)$ as $s\to\infty$ that
is independent of $n$.

\item The reduced model is the leading asymptotic term of the main regime.
\Crefrange{prop:decomp}{prop:var}, \cref{lem:cond} and
\crefrange{cor:bounds}{cor:srvar} are rigorous, while \cref{eq:reduced} discards the
$O(\varepsilon)$ term of \cref{lem:cond}, whose inequality degenerates when $2^\xi$
lies far below the data scale (\cref{lem:cond}). It matches the simulations to a
few per cent in the experimental main regime but is untested in the
dead-zone-dominated one, and the $L_m$ scaling law $\rho\propto4^{-L_m}$
(\cref{eq:lmscaling}) likewise fails near the peak of $|\delta(s)|$: the transition
region lacks a rigorous description, and the test of (A2) in \cref{sec:cond} is
binned, hence numerical rather than analytic.

\item The format model is idealized, and two observations are informal.
\Cref{eq:format} assumes the shared exponent to be the rounded logarithmic scale of
the largest magnitude in the block, round-to-nearest rounding, and no overflow; we
tested four rounding rules but not real hardware, and accumulator width, sticky-bit
handling and normalization in tensor cores could all change the conclusions.
Separately, in initial experiments the growth order of the dot-product error under
a shared exponent was worse than under an elementwise exponent, but no
same-standard comparison was carried out, and the literature search comprised two
focused passes.
\end{enumerate}

\section{Conclusions}
\label{sec:conclusions}

This paper re-examines a statistical assumption carried in block floating point
error analyses since 1996: that within-block quantization errors are zero-mean and
uncorrelated. The covariance identity
$\operatorname{Cov}(e_i,e_j)=\mathbb{E}[\operatorname{Cov}(e_i,e_j\mid\xi)]
+\operatorname{Var}(m(\xi))$ needs no extra assumptions, conditional
exchangeability being automatic, and Cauchy--Schwarz gives
$0\le\mathcal I_n\le n$, so any asymptotic log--log slope is at most $1$;
\cref{prop:decay} sharpens this to $\mathcal I_n=o(n)$ for unbounded support with
$\mathbb{E}[x_1^2]<\infty$. The variance law separates two independent defects of
the classical model: $\delta\not\equiv0$ creates correlation, and
$\gamma\not\equiv1/12$ biases the variance formula. Under stochastic rounding
both terms vanish, so $\mathcal I_n\equiv1$.

The numerical results support the same picture. For symmetric distributions, (H1)
and (H3) hold in the measured range; for asymmetric ones, a direct decomposition
given $\xi$ shows that the within-block correlation comes almost entirely from
$\operatorname{Var}(m(\xi))$. With no fitted parameter, the reduced model
reproduces measured $\mathcal I_n$ within $4.5\%$ over $36$ round-to-nearest
configurations, predicts the dead-zone scaling $\rho\propto4^{-L_m}$ and the ReLU
variance ratio $0.506$, and yields $1.000$ under stochastic rounding by the same
atom argument.

The asymptotic classification is clear: bounded support forces $\mathcal I_n\to1$
at an exponential rate (\cref{prop:bound}), while unbounded support produces an
inflation that first rises and then falls, \cref{prop:decay} ruling out asymptotic
linear accumulation for any distribution in this class with finite second moment.
Large-$n$ measurements and exact quadrature of the reduced model agree on this
fall-off (\cref{sec:inflbig}).

The practical message is narrow but concrete: current hardware need not already be
inaccurate at the MX block length $n=32$, the effect size depending strongly on
the input distribution. What does fail, once the data are asymmetric, is the
independence premise used to justify $\sqrt n$ accumulation in probabilistic
error analyses. BFP models applied to neural-network activations should therefore
be corrected, and the scale of that failure is controlled: it appears at
moderate block lengths and then decays, rather than growing without bound.

\clearpage
\appendix

\section{Data availability}
\label{app:full}

The complete tables behind \cref{sec:results} — the full grid of
$\mathcal I_n$ over six distributions, two mantissa widths and seven block
lengths, and the stochastic-rounding comparison at $n=32$, $512$ and $4096$ —
are distributed with the reproduction package described in
\cref{sec:methods}, together with the archived terminal output of every batch of
runs. The tables printed in \cref{sec:results} are the subsets that the
argument in this paper uses; the only values quoted in the text but not tabulated
there are the $n=32$ and $n=512$ entries of \cref{tab:srfull} and the
$z$ values of \cref{fig:biasvar} (top), all of which appear in the package.

\section*{Acknowledgments}
The author thanks the College of Computer Science and Technology, Nanjing
University of Aeronautics and Astronautics, for the academic environment in
which this work was carried out.

The author used machine-translation and large-language-model assistance for
drafting and language editing of the English text. All mathematical content,
all numerical results and their reproduction scripts, and all bibliographic
entries were checked against the primary sources by the author, who takes full
responsibility for the content of the manuscript.

\bibliographystyle{siamplain}
{\small\bibliography{references}}
\end{document}
